% !TEX program = pdflatex
% Aura: Гібридний постквантовий протокол захищеного обміну повідомленнями
% Єдина українська версія статті

\documentclass[11pt,a4paper]{article}
% Преамбула єдиної української статті

% ─── Геометрія ───────────────────────────────────────────────
\usepackage[margin=2.5cm]{geometry}

% ─── Кодування та шрифти (pdflatex + T2A для кирилиці) ──────
\usepackage[T2A,T1]{fontenc}
\usepackage[utf8]{inputenc}
\usepackage[english,ukrainian]{babel}
% Latin Modern does not provide T2A Cyrillic shapes in many pdfLaTeX installs;
% the default LH/CM family covers Ukrainian text without font-shape fallbacks.
% T2A Cyrillic fonts are often bitmap-backed in minimal TeX installs; disable
% expansion so pdfLaTeX does not fail when scalable Cyrillic fonts are absent.
\usepackage[protrusion=true,expansion=false]{microtype}

% ─── Математика ─────────────────────────────────────────────
\usepackage{amsmath,amssymb,amsthm}

% ─── Таблиці та рисунки ─────────────────────────────────────
\usepackage{booktabs}
\usepackage{tabularx}
\usepackage{multirow}
\usepackage{xcolor}
\usepackage{graphicx}
\usepackage{float}
\usepackage{caption}
\usepackage{subcaption}
\newcolumntype{Y}{>{\raggedright\arraybackslash}X}
\setlength{\emergencystretch}{3em}

% ─── Алгоритми ───────────────────────────────────────────────
\usepackage[ruled,linesnumbered,noend]{algorithm2e}
\SetKwComment{Comment}{$\triangleright$~}{}
\SetAlgorithmName{Алгоритм}{алгоритм}{Список алгоритмів}

% ─── Гіперпосилання ─────────────────────────────────────────
\usepackage[colorlinks,citecolor=blue!70!black,linkcolor=blue!70!black,urlcolor=blue!50!black]{hyperref}
\usepackage[ukrainian,capitalise,noabbrev]{cleveref}

% ─── Лістинги коду ──────────────────────────────────────────
\usepackage{listings}
\lstset{
  basicstyle=\ttfamily\small,
  breaklines=true,
  frame=single,
  columns=fullflexible,
}

% ─── Теоремні середовища (українською) ──────────────────────
\newtheorem{theorem}{Теорема}[section]
\newtheorem{lemma}[theorem]{Лема}
\newtheorem{corollary}[theorem]{Наслідок}
\newtheorem{proposition}[theorem]{Твердження}
\theoremstyle{definition}
\newtheorem{definition}[theorem]{Означення}
\newtheorem{remark}[theorem]{Зауваження}
\renewcommand{\proofname}{Доведення}

% ─── Нотація: загальні криптографічні позначення ────────────
\newcommand{\secpar}{\ensuremath{\lambda}}
\newcommand{\negl}{\ensuremath{\mathsf{negl}}}
\newcommand{\ppt}{\ensuremath{\mathsf{PPT}}}
\newcommand{\adv}{\ensuremath{\mathbf{Adv}}}
\newcommand{\advA}{\ensuremath{\mathcal{A}}}
\newcommand{\advB}{\ensuremath{\mathcal{B}}}
\newcommand{\advS}{\ensuremath{\mathcal{S}}}
\newcommand{\challenger}{\ensuremath{\mathcal{C}}}
\newcommand{\oracle}{\ensuremath{\mathcal{O}}}

\newcommand{\HKDF}{\textsf{HKDF}}
\newcommand{\HMAC}{\textsf{HMAC}}
\newcommand{\SHA}{\textsf{SHA-256}}
\newcommand{\AEAD}{\textsf{AES-256-GCM-SIV}}
\renewcommand{\DH}{\textsf{X25519}}
\newcommand{\KEM}{\textsf{Kyber-768}}
\newcommand{\KeyGen}{\textsf{KeyGen}}
\newcommand{\Encaps}{\textsf{Encaps}}
\newcommand{\Decaps}{\textsf{Decaps}}
\newcommand{\Enc}{\textsf{Enc}}
\newcommand{\Dec}{\textsf{Dec}}

\newcommand{\GapCDH}{\textsf{Gap-CDH}}
\newcommand{\DDH}{\textsf{DDH}}
\newcommand{\INDCCA}{\textsf{IND-CCA2}}
\newcommand{\INDCPA}{\textsf{IND-CPA}}
\newcommand{\INTCTXT}{\textsf{INT-CTXT}}
\newcommand{\MRAE}{\textsf{MRAE}}
\newcommand{\PRF}{\textsf{PRF}}
\newcommand{\dualPRF}{\textsf{dual-PRF}}
\newcommand{\ROM}{\textsf{ROM}}
\newcommand{\eCK}{\textsf{eCK}}

\newcommand{\concat}{\ensuremath{\,\|\,}}
\newcommand{\getsr}{\ensuremath{\stackrel{\$}{\gets}}}

% ─── Нотація: ключі та епохи протоколу Aura ─────────────────
\newcommand{\rk}{\ensuremath{\mathit{rk}}}
\newcommand{\ck}{\ensuremath{\mathit{ck}}}
\newcommand{\mk}{\ensuremath{\mathit{mk}}}
\newcommand{\mdk}{\ensuremath{\mathit{mdk}}}
\newcommand{\epochv}{\ensuremath{\mathit{epoch}}}
\newcommand{\IK}{\ensuremath{\mathit{IK}}}
\newcommand{\SPK}{\ensuremath{\mathit{SPK}}}
\newcommand{\EK}{\ensuremath{\mathit{EK}}}
\newcommand{\OPK}{\ensuremath{\mathit{OPK}}}

% ─── Зручні скорочення для тексту ───────────────────────────
\newcommand{\Aura}{\mbox{\normalfont\scshape Aura}}
\newcommand{\Signal}{\mbox{\normalfont\scshape Signal}}
\newcommand{\PQXDH}{\textsf{PQXDH}}
\newcommand{\SPQR}{\textsf{SPQR}}
\newcommand{\PQThree}{\textsf{PQ3}}
\newcommand{\MLS}{\textsf{MLS}}

% ─── Метадані заголовку ─────────────────────────────────────
\newcommand{\paperauthor}{Олександр Мельниченко\\\texttt{oleksandr@aura.io}}

\title{\textbf{Aura: гібридний постквантовий протокол захищеного обміну\\
повідомленнями з KEM-ротацією на кроках ратчета}}
\author{\paperauthor}
\date{\today}

\begin{document}
\maketitle

\begin{abstract}
\noindent
Подано \Aura{} — гібридний постквантовий протокол захищеного обміну
повідомленнями, що поєднує X25519 і ML-KEM-768 у початковому
X3DH-подібному рукостисканні та в подальших напрямкових кроках ратчета.
Кожна зміна напрямку повідомлень деривує новий кореневий ключ зі свіжого
\DH{}-секрету та KEM-внеску. У моделі компрометації, розглянутій у
статті, це дає класичне PCS-відновлення після одного чесного
ratchet-кроку та гібридне відновлення після двох напрямкових кроків.

Конструкція прив'язує кожний новий KEM-публічний ключ до HKDF-транскрипту
та відділяє шифрування корисного навантаження від шифрування метаданих
через незалежний ключ метаданих епохи. Порівняно із \Signal{}
\SPQR{}/Triple Ratchet та Apple \PQThree{} внесок цієї статті обмежено
відкритим Rust-артефактом, у якому KEM-key binding, metadata-key rotation
і proof-to-code зв'язки можна перевіряти в одній реалізації.

Для протоколу сформульовано модель безпеки в стилі \eCK{} і наведено
шість тверджень безпеки: стійкість гібридного KEM-комбінатора,
\eCK{}-AKE-стійкість, пряму секретність, постскомпрометаційне
відновлення з доведеною асиметрією «1~класичний крок / 2~гібридні
кроки», конфіденційність та цілісність повідомлень, а також стійкість
до повторів. До технічних результатів роботи належать доповнений HKDF з
прив'язуванням нового KEM-ключа до параметра $\mathit{info}$,
дворівневе AEAD з ротаційним ключем метаданих, формалізація
межі KEM-асиметрії PCS для цієї конструкції та відображення доказових
інваріантів на Rust-реалізацію. Символьну перевірку виконано в
Tamarin~Prover (10/10 лем) і ProVerif (3/6 зобов'язань для handshake/message-path
доведено; Q3/Q5/Q6 не враховано як машинно-доведені результати). Реалізація мовою
Rust містить близько 35{,}9~тис. рядків у \texttt{src/}; повний репозиторний
зріз тестів містить 750~сценаріїв у стандартній конфігурації та
889~сценаріїв із прапорцем \texttt{ffi}.
\end{abstract}

\smallskip
\noindent\textbf{Ключові слова:} постквантова криптографія, захищений
обмін повідомленнями, Double Ratchet, X3DH, ML-KEM, HKDF,
AES-GCM-SIV, MRAE, Tamarin Prover, ProVerif, Rust.

\medskip
\noindent\textit{English abstract.}
\Aura{} is a hybrid post-quantum messaging protocol that combines X25519
ECDH with ML-KEM-768 in the initial handshake and in subsequent
bidirectional ratchet steps. The post-quantum material is treated as part
of the state-healing mechanism, not as a one-time X3DH extension. The
protocol uses per-direction KEM rotation, HKDF binding of fresh KEM public
keys, two-layer AEAD with a rotating metadata key, and a formal
1-step classical / 2-step hybrid PCS recovery result. The analysis
includes an \eCK{}-style model, six security claims, symbolic checks in
Tamarin Prover and ProVerif, and a Rust implementation whose full repository
test matrix contains 750/889 scenarios depending on the \texttt{ffi} feature.

\tableofcontents
\newpage


% ═════════════════════════════════════════════════════════════
% 1. ВСТУП
% ═════════════════════════════════════════════════════════════
\section{Вступ}\label{sec:intro}

\subsection{Актуальність дослідження}\label{sec:relevance}

Протокол \Signal{}~\cite{signal-x3dh,cohn-gordon2020} є широко застосованою
основою наскрізного шифрування у месенджерах: на ньому ґрунтуються
Signal~Messenger, WhatsApp, Google Messages, RCS. Обмін ключами X3DH та
алгоритм подвійного ратчета (Double Ratchet) забезпечують пряму секретність
(forward secrecy, FS) та постскомпрометаційну безпеку (post-compromise
security, PCS) за класичних припущень про ECDH на Curve25519.

Поява криптографічно-релевантних квантових комп'ютерів (CRQC) ставить під
загрозу всі протоколи, що покладаються виключно на класичні припущення про
складність задач дискретного логарифмування та факторизації. Алгоритм
Шора~\cite{shor1997} ефективно розв'язує обидві задачі на достатньо великому
квантовому комп'ютері. Хоча NIST стандартизував гратковий механізм
інкапсуляції ML-KEM (FIPS~203~\cite{fips203}), пряма заміна класичних
примітивів на постквантові у протоколах двостороннього зв'язку наражається
на чотири системні проблеми, які формулюємо нижче як \emph{виклики}.

\paragraph{Виклик~1. Атаки типу «зберегти зараз — розшифрувати пізніше»
(\emph{harvest-now-decrypt-later}).}
Зловмисник, що сьогодні записує зашифрований ECDH-трафік, може згодом
відновити відповідні сесійні секрети після появи CRQC. Пряма секретність на
основі ефемерного ECDH не захищає саме від цієї моделі загрози, бо алгоритм
Шора ретроспективно розв'язує DLP.

\paragraph{Виклик~2. Гранулярність інтеграції KEM.}
Специфікація PQXDH~\cite{signal-pqxdh} додає єдину інкапсуляцію Kyber до
початкового рукостискання, але \emph{не} включає постквантовий ключовий
матеріал у наступні кроки ратчета. Отже, постскомпрометаційне відновлення
після handshake спирається на класичний DH, а квантовостійка PCS-властивість
для наступних епох не випливає з моделі PQXDH.

\paragraph{Виклик~3. Асиметрія KEM-відновлення.}
На відміну від DH, де обидві сторони симетрично надають ефемерні ключі,
KEM вимагає, щоб одна сторона першою опублікувала публічний ключ, а інша
лише потім могла інкапсулювати. Це породжує внутрішню асиметрію
гібридних ратчет-протоколів: класичне PCS-відновлення відбувається
за один крок ратчета, але повне гібридне (квантово-стійке) відновлення
вимагає двох кроків. Цю асиметрію необхідно формалізувати в моделі
безпеки, а не вважати недоліком реалізації.

\paragraph{Виклик~4. Витік метаданих.}
Стандартні реалізації Double Ratchet передають публічні ключі ратчета,
лічильники епохи та індекси повідомлень у відкритому вигляді в заголовку
конверта. Навіть за умови шифрування вмісту такі метадані розкривають
комунікаційні шаблони, порядок повідомлень та життєвий цикл сесії
мережевому спостерігачу.

\medskip
У цій роботі ці чотири питання розглядаються разом. PQXDH є коректною
точкою порівняння для початкового handshake, але не для подальшої
постквантової еволюції стану. Sealed~Sender у \Signal{} працює на
іншому рівні — приховує відправника від сервісу, але не змінює формат
Double Ratchet-заголовка. Для \Aura{} важливий саме протокольний шар,
де KEM-матеріал, metadata-key і replay/sealed-state інваріанти мають
бути узгоджені в одному стані сесії.

\subsection{Мета та завдання дослідження}\label{sec:goals}

\textbf{Мета.} Сконструювати та формально проаналізувати протокол
наскрізного шифрування для двостороннього обміну повідомленнями, що
одночасно: \emph{(а)}~аналізує гібридну конфіденційність упродовж
ратчетних епох у межах визначеної моделі загроз; \emph{(б)}~зберігає
пряму секретність та постскомпрометаційну безпеку з явно вказаними
умовами; \emph{(в)}~захищає метадані протокольного конверта без
заявок на анонімність чи traffic-analysis resistance; \emph{(г)}~допускає
машинно-перевірену символьну верифікацію ключових властивостей.

\textbf{Завдання:}
\begin{enumerate}
  \item Систематизувати та формалізувати чотири виклики постквантового
    дизайну месенджерних протоколів.
  \item Запропонувати гібридне X3DH-рукостискання, у якому
    спільний секрет утворюється як комбінатор класичного DH-секрету
    та постквантового KEM-секрету.
  \item Розробити гібридний Double Ratchet, що на кожній зміні напрямку
    породжує нову пару \DH{}- та \KEM{}-ключів, із доповненим (augmented)
    HKDF, який криптографічно прив'язує новий KEM-ключ до похідного
    кореневого ключа.
  \item Сконструювати дворівневе AEAD-шифрування з ротаційним ключем
    метаданих для одночасного захисту корисного навантаження та
    заголовка конверта, спираючись на MRAE-схему \AEAD{}.
  \item Запропонувати механізм запечатаного (sealed) збереження стану
    сесії з контролем відкату за монотонним зовнішнім лічильником.
  \item Розробити формальну модель безпеки в стилі \eCK{} та довести
    шість теорем безпеки.
  \item Виконати машинно-перевірену символьну верифікацію
    (Tamarin~Prover, ProVerif) ключових властивостей.
  \item Розробити повну реалізацію мовою Rust і провести експериментальну
    оцінку.
\end{enumerate}

\subsection{Об'єкт та предмет дослідження}\label{sec:object}

\textbf{Об'єкт дослідження} — криптографічні протоколи захищеного
двостороннього обміну повідомленнями зі станом, що підтримують пряму
секретність та постскомпрометаційну безпеку у моделях класичних та
квантових опонентів.

\textbf{Предмет дослідження} — методи та архітектурні рішення інтеграції
постквантових механізмів інкапсуляції ключів у схеми X3DH/Double~Ratchet
із явними умовами прямої секретності та PCS упродовж ратчетних епох,
у тому числі: \emph{(а)}~гібридні комбінатори ключового матеріалу;
\emph{(б)}~KEM-ротація на кроках ратчета; \emph{(в)}~доповнене
прив'язування KEM-ключа до деривації ратчета через параметр $\mathit{info}$
схеми HKDF; \emph{(г)}~дворівневе MRAE-шифрування з ротаційним ключем
метаданих.

\subsection{Наукова новизна одержаних результатів}\label{sec:novelty}

Наукова новизна роботи зосереджена на становій еволюції сесії. У
\PQXDH{} KEM-секрет входить лише до початкового ключового матеріалу;
після цього PCS знову спирається на класичний DH. У \Aura{} KEM-пара є
частиною напрямкового ratchet-переходу, тому постквантовий внесок
регулярно повертається в кореневий ключ. Порівняно з \SPQR{} та
\PQThree{} ця робота фіксує інший профіль: відкриту Rust-конструкцію з
proof-to-code відображенням, явним HKDF-зв'язуванням нового KEM-ключа
та окремою ротацією ключа метаданих.

Це формулює окрему наукову тезу: для двостороннього месенджера
постквантова безпека не зводиться до постквантового X3DH. Вона потребує
конструкції, яка одночасно задає: \emph{(а)}~як KEM-внесок оновлюється
після компрометації стану; \emph{(б)}~як новий KEM-ключ зв'язується з
похідним станом без підпису на кожному ratchet-кроці; \emph{(в)}~яка мінімальна затримка
гібридного PCS-відновлення випливає з односторонності KEM; \emph{(г)}~як
доказові інваріанти зберігаються в реалізації.

\medskip
\noindent\textbf{Основні нові результати:}

\begin{enumerate}
  \item[\textbf{Н1.}] \textbf{Напрямковий гібридний Double Ratchet з KEM-ротацією.}
    Запропоновано та реалізовано конструкцію ратчета,
    у якій \emph{кожна} зміна напрямку повідомлень викликає
    гібридне оновлення на основі ефемерного \DH{}-секрету та
    свіжої пари \KEM{}-ключів. На відміну від PQXDH, де постквантовий
    ключовий матеріал входить лише до початкового рукостискання, у \Aura{}
    він входить до кореневих ключів, що виникають на ратчетних змінах
    напрямку. На відміну від sparse/periodic підходів
    \SPQR{} та \PQThree{}, тут досліджується компактний per-ratchet
    профіль з окремим proof-to-code відображенням і явним
    твердженням про затримку PCS-відновлення.

  \item[\textbf{Н2.}] \textbf{Доповнений (augmented) HKDF для
    перевірюваного зв'язування KEM-ключа з похідним станом.}
    Запропоновано прив'язувати новий публічний \KEM{}-ключ до
    параметра $\mathit{info}$ функції деривації ратчета через конкатенацію
    $\mathit{info} := \texttt{Aura-Hybrid-Ratchet} \concat \mathit{pk\_new}$.
    Якщо опонент модифікує $\mathit{pk\_new}$ у транзиті, обидві сторони
    одержать різні $\mathit{info}$ і, як наслідок, різні кореневі ключі; це
    призводить до невдалого AEAD-розшифрування наступного повідомлення.
    У такий спосіб досягається перевірюване зв'язування KEM-ключа з
    похідним станом без застосування додаткових цифрових підписів і без
    додаткового round-trip між сторонами: підміна призводить до
    розходження ключів і наступного AEAD reject. HKDF використовується не
    лише як примітив деривації, а й як механізм зв'язування ратчетного
    KEM-ключа з наступним кореневим станом.

  \item[\textbf{Н3.}] \textbf{Дворівневе AEAD з ротаційним ключем
    метаданих.}
    Запропоновано конструкцію, в якій метадані конверта
    (тип повідомлення, індекс, ідентифікатор кореляції) шифруються
    окремим ключем $\mdk_{\epochv}$, унікальним для кожної ратчетної
    епохи, що деривується одночасно з кореневим та ланцюговим ключами
    із виходу HKDF загального обсягу $|\rk| + |\ck| + |\mdk|$. Корисне
    навантаження шифрується незалежним ключем повідомлення, причому
    обидва рівні використовують MRAE-схему \AEAD{}, стійку до повторного
    використання nonce (RFC~8452). Це забезпечує \emph{(а)}~захист
    метаданих від спостерігача мережі; \emph{(б)}~додатковий рівень захисту
    проти випадкового повторного nonce у разі катастрофічної помилки
    реалізації.

  \item[\textbf{Н4.}] \textbf{Формалізація асиметрії PCS-відновлення.}
    Явно сформульовано і доведено для \Aura{} межу відновлення
    «1~крок для класичної PCS / 2~кроки для гібридної PCS» у
    консервативній моделі компрометації обох кінців. Додатковий
    гібридний крок трактується як наслідок напрямковості KEM
    (спочатку треба опублікувати новий KEM-ключ, потім отримати
    інкапсуляцію до нього), а не як артефакт реалізації.

  \item[\textbf{Н5.}] \textbf{Комплексна формальна верифікація гібридного
    постквантового месенджера.}
    Для запропонованого двостороннього протоколу подано одночасно:
    \emph{(а)}~шість теорем безпеки і вісім допоміжних лем із
    доведеннями методом редукцій до стандартних припущень
    (\GapCDH{}, \INDCCA{} для ML-KEM, \dualPRF{}, \MRAE{});
    \emph{(б)}~машинно-перевірену верифікацію в Tamarin~Prover
    (10 з 10 лем); \emph{(в)}~машинно-перевірену верифікацію в
    ProVerif (3 з 6 зобов'язань, з аналізом незавершення, меж запиту та
    очікуваної хибності stress-запиту для трьох відкритих пунктів); \emph{(г)}~Rust-тести, benchmark-и
    та FFI-межу, прив'язані до тієї самої протокольної машини станів.
\end{enumerate}

\def\NoveltyPCSRef{\cref{thm:pcs}}
\def\NoveltyConfRef{\cref{thm:conf-int}}
\def\NoveltyDiscRef{\S\ref{sec:discussion}}
\def\NoveltyFormalRefs{\S\ref{sec:formal}, \S\ref{sec:code-proof-map},
\S\ref{sec:negative-scenarios}}

\begin{table}[H]
\centering
\caption{Матриця наукової новизни та її підтвердження.}
\label{tab:novelty-matrix}
\small
\begin{tabularx}{\textwidth}{@{}p{0.15\textwidth}YYY@{}}
\toprule
Результат & Що є новим & Відмінність від наявних підходів & Де підтверджено \\
\midrule
\textbf{Н1} &
KEM-ротація на кожному напрямковому ratchet-кроці &
PQXDH додає KEM лише до handshake; \SPQR{} та \PQThree{} мають
sparse/periodic PQ-пере-ключення; \Aura{} фіксує per-ratchet профіль
оновлення кореневого ключа &
\S\ref{sec:ratchet}, \NoveltyPCSRef, тести PCS \\
\addlinespace
\textbf{Н2} &
Прив'язування $\mathit{pk\_new}$ до $\mathit{info}$ у HKDF &
Новий KEM-ключ не підписується окремо, а автентифікується через
розбіжність похідних ключів при підміні &
\S\ref{sec:augmented}, Tamarin-модель ратчета, тести підміни \\
\addlinespace
\textbf{Н3} &
Окремий ротаційний metadata-key на кожну епоху &
Метадані envelope не лишаються відкритим побічним каналом ратчета &
\S\ref{sec:two-layer}, \NoveltyConfRef, metadata-тести \\
\addlinespace
\textbf{Н4} &
Мінімальна асиметрія PCS: 1 класичний / 2 гібридні кроки &
Затримка в два кроки пояснена як наслідок KEM, а не як недолік
реалізації &
\NoveltyPCSRef, \NoveltyDiscRef \\
\addlinespace
\textbf{Н5} &
Єдиний ланцюг: модель, теореми, Tamarin/ProVerif, Rust-тести &
Доказові твердження пов'язані з інваріантами стану, replay і FFI &
\NoveltyFormalRefs \\
\bottomrule
\end{tabularx}
\end{table}

\subsection{Практичне значення}\label{sec:practical}

Запропонований протокол \Aura{} реалізовано мовою Rust (близько
35{,}9~тис. рядків у \texttt{src/} та 30{,}1~тис. рядків тестового коду;
750~тестових сценаріїв у стандартній конфігурації та 889~сценаріїв
із прапорцем \texttt{ffi}) із FFI-обгорткою для крос-платформного
використання (iOS/Swift, Android, Windows/.NET, серверна сторона на Rust).
Реалізація відкрита та проходить безперервну інтеграцію з
автоматичним прогоном KAT-векторів (RFC~5869, RFC~7748, RFC~8452),
property-based-тестів та конкурентних стрес-тестів. Латентність
рукостискання $\sim$1{,}1~мс і пропускна здатність у режимі
бурсту $\sim$15~мкс/повідомлення на Apple~M-series показують, що у
виміряному профілі витрати сумісні з асинхронним текстовим сценарієм.

\subsection{Структура статті}\label{sec:structure}

У розділі~\ref{sec:related} оглянуто споріднені роботи. У
розділі~\ref{sec:prelim} введено нотацію та означення криптографічних
примітивів. У розділі~\ref{sec:protocol} викладено архітектуру
протоколу: пре-ключовий бандл (\S\ref{sec:bundle}), гібридний X3DH
(\S\ref{sec:handshake}), гібридний Double Ratchet (\S\ref{sec:ratchet})
із доповненим HKDF (\S\ref{sec:augmented}), дворівневе AEAD
(\S\ref{sec:two-layer}), захист від повторів (\S\ref{sec:replay}) та
запечатане збереження стану (\S\ref{sec:sealed}).
Далі у розділах~\ref{sec:model}--\ref{sec:implementation} та
\ref{sec:evaluation} подано модель безпеки, доведення, машинну
верифікацію, реалізацію та експериментальну оцінку.

% ═════════════════════════════════════════════════════════════
% 2. СПОРІДНЕНІ РОБОТИ
% ═════════════════════════════════════════════════════════════
\section{Споріднені роботи}\label{sec:related}

\paragraph{Сім'я протоколів \Signal{}.}
Протокол \Signal{}, формалізований Cohn-Gordon et al.~\cite{cohn-gordon2020}
та подальшими аналізами X3DH/ратчет-сімейства~\cite{brendel2020,hashimoto2021},
досягає прямої секретності та постскомпрометаційної безпеки шляхом поєднання
X3DH-обміну ключами~\cite{signal-x3dh} та алгоритму Double Ratchet. Протокол
зазнав ретельного формального аналізу як у символьній, так і в обчислювальній
моделях. Класичний \Signal{} ґрунтується виключно на припущенні \GapCDH{}
для Curve25519, тому його DH-похідні сесійні ключі не захищають записаний
сьогодні трафік від майбутнього опонента, здатного розв'язувати ECDLP.

\paragraph{Постквантові розширення \Signal{}.}
Специфікація PQXDH~\cite{signal-pqxdh} додає одну інкапсуляцію Kyber-1024
до X3DH-рукостискання, чим адресує атаки «зберегти зараз --- розшифрувати
пізніше» для початкового handshake. Проте PQXDH \emph{не} вносить
постквантовий ключовий матеріал у наступні ратчет-кроки: подальше
постскомпрометаційне відновлення після handshake спирається на класичний
DH, тож квантовостійка PCS-властивість для наступних епох не випливає з
моделі PQXDH.

До порівняння також слід включати \SPQR{} та Triple
Ratchet~\cite{signal-spqr,signal-double-ratchet}: постквантовий
Sparse Continuous Key Agreement працює поруч із класичним Double
Ratchet, а ключі двох механізмів змішуються у гібридний ключ
повідомлення. Базовий SCKA-компонент Signal описано як ML-KEM
Braid~\cite{mlkem-braid}. Отже, \Aura{} потрібно порівнювати не лише
з handshake-only схемами, а й з протоколами, де PQ-матеріал бере участь
в еволюції стану.

\paragraph{Промислові постквантові месенджери.}
Apple \PQThree{}~\cite{apple-pq3,stebila-pq3,linker-pq3} також
переносить постквантовий внесок з початкового обміну у подальше
оновлення ключів: протокол має гібридне встановлення сесії, періодичне
постквантове пере-ключення та незалежний формальний аналіз. Тому
\Aura{} не слід подавати як першу у світі ідею ``PQ у ратчеті''.
Новизна цієї роботи вужча й точніша: відкрита Rust-реалізація
двостороннього протоколу, де KEM-матеріал ротують на межах ратчета,
новий KEM-ключ зв'язують через HKDF, метадані мають окремий ротаційний
ключ, а доказові інваріанти зіставлено з кодом.

\paragraph{Інші протоколи з гібридним обміном.}
TLS~1.3 + Kyber (експериментальний CECPQ2~\cite{cecpq2}, NIST PQC TLS
hybrid drafts) застосовують гібрид у TLS-рукостисканні, але не мають
ратчет-структури; це продовжує лінію робіт із постквантового
обміну ключами для Інтернету~\cite{stebila2020} та формального
аналізу TLS~1.3~\cite{dowling2020}. Протокол MLS~\cite{rfc9420} орієнтований на групове
шифрування і забезпечує епохальну FS/PCS для груп, але базовий RFC
9420 не є двостороннім Double-Ratchet-подібним протоколом і не
стандартизує постквантовий внесок за замовчуванням. У цій статті
\MLS{} використовується як референс для групової еволюції стану, а
не як прямий конкурент двосторонньому месенджерному каналу.

\paragraph{Шифрування метаданих месенджерів.}
Sealed~Sender~\cite{sealed-sender} приховує відправника від сервера,
але не приховує заголовок Double Ratchet від мережного спостерігача.
Підхід \Aura{} ортогональний: ми шифруємо заголовок конверта, а не
переписуємо систему доставки. Sphinx~\cite{sphinx}, Loopix~\cite{loopix}
працюють на рівні mix-мереж і не вирішують задачу заголовка
end-to-end-протоколу.

\paragraph{Дослідження ратчет-схем.}
Alwen, Coretti, Dodis~\cite{acd2019} запровадили теоретичну рамку для
аналізу безперервного обміну ключами (Continuous Key Agreement) та
показали межі досяжного PCS. Bienstock et al.~\cite{bienstock2020}
розглянули постквантові ратчети у моделі ідеальних KEM, однак не
дослідили конкретні архітектурні прийоми (як-от доповнений HKDF)
для зв'язування нового KEM-ключа з похідним станом. Для групових
ратчетів окремо важлива ціна конкурентності~\cite{bienstock2022},
що пояснює, чому в цій роботі двосторонній канал і групові схеми
розглянуто як різні об'єкти аналізу.

% ═════════════════════════════════════════════════════════════
% 3. ПЕРЕДУМОВИ
% ═════════════════════════════════════════════════════════════
\section{Криптографічні передумови}\label{sec:prelim}

\subsection{Нотація}\label{sec:notation}

Для $\secpar \in \mathbb{N}$ пишемо $\{0,1\}^{\secpar}$ для множини бітових
рядків довжини $\secpar$ та $\{0,1\}^{*}$ для множини бітових рядків
скінченної довжини. Конкатенацію двох рядків $a, b$ позначаємо
$a \concat b$. Випадковий рівномірний вибір елемента $x$ зі скінченної
множини $S$ позначаємо $x \getsr S$.

Функцію називаємо \emph{знехтуваною} (negligible), якщо для будь-якого
многочлена $p$ існує $\secpar_0$ таке, що $f(\secpar) < 1/p(\secpar)$ для
всіх $\secpar > \secpar_0$. Множину знехтуваних функцій позначаємо
$\negl(\secpar)$. Усі опоненти ($\advA$, $\advB$, $\advS$) розглядаються
як ймовірнісні поліноміальні машини Тюрінга (\ppt).

\subsection{Криптографічні примітиви}\label{sec:primitives}

\begin{definition}[Обмін ключами Діффі--Геллмана на X25519]
\label{def:dh}
У доведенні X25519 ідеалізовано як DH у великій prime-order підгрупі після
clamping та all-zero/small-order rejection; це не буквальна модель усієї
Curve25519 з кофактором 8. Функцію \DH{} записуємо як: при секретному
скалярі $a \in \mathbb{Z}_q$ та публічному ключі $B = g^{b}$
обчислюється $\DH(a, B) = B^{a} = g^{ab}$. Стандарт
RFC~7748~\cite{rfc7748} задає X25519-кодування, clamping і вимоги до
обробки малопорядкових точок.
\end{definition}

\begin{definition}[\GapCDH{}]
\label{def:gap-cdh}
Для \DH{}-групи $\mathbb{G}$ задача \GapCDH{} полягає в обчисленні
$g^{ab}$ за заданими $(g,g^a,g^b)$ за наявності оракула перевірки
\DDH{}. Перевагу опонента позначаємо
\begin{equation}\label{eq:gap-cdh}
  \adv^{\GapCDH}_{\DH}(\advA)
  =
  \Pr[\advA(g,g^a,g^b;\mathcal{O}_{\DDH}) = g^{ab}].
\end{equation}
\end{definition}

\begin{definition}[KEM]
\label{def:kem}
KEM — це трійка $(\KeyGen, \Encaps, \Decaps)$:
\begin{itemize}
  \item $\KeyGen(1^{\secpar}) \to (\mathit{pk}, \mathit{sk})$ — генератор
    пари ключів;
  \item $\Encaps(\mathit{pk}) \to (c, K)$ — інкапсуляція: видає шифротекст
    $c$ та секретний ключ $K \in \{0,1\}^{\secpar}$;
  \item $\Decaps(\mathit{sk}, c) \to K$ або $\bot$ — декапсуляція.
\end{itemize}
Кажемо, що KEM є \INDCCA{}-стійким, якщо перевага будь-якого \ppt{}
опонента у грі IND-CCA2 знехтувано мала:
\begin{equation}\label{eq:indcca}
  \adv^{\INDCCA}_{\KEM}(\advA)
  =
  \left|\Pr[\mathsf{Exp}^{\INDCCA}_{\KEM}(\advA)=1]-\frac12\right|.
\end{equation}
Використовуваний у \Aura{}
ML-KEM-768 (FIPS~203~\cite{fips203}, NIST security category~3) задовольняє цю
вимогу при стандартних припущеннях про складність задачі \textsf{Module-LWE}.
\end{definition}

\begin{definition}[HKDF як подвійна PRF]
\label{def:hkdf}
HKDF~\cite{rfc5869} — це конструкція extract-then-expand на основі HMAC:
\begin{align}
  \mathit{prk} &\gets \mathsf{Extract}(\mathit{salt}, \mathit{ikm}) = \HMAC_{\mathit{salt}}(\mathit{ikm}), \label{eq:hkdf-extract}\\
  \mathit{okm} &\gets \mathsf{Expand}(\mathit{prk}, \mathit{info}, L). \label{eq:hkdf-expand}
\end{align}
У моделі випадкового оракула HKDF поводиться як подвійна PRF
(\dualPRF{}): при фіксованому невідомому $\mathit{salt}$ й випадковому
$\mathit{ikm}$ вихід є псевдовипадковим, та симетрично — при
випадковому $\mathit{salt}$ і фіксованому $\mathit{ikm}$ вихід також
псевдовипадковий~\cite{krawczyk2010}.
\end{definition}

\begin{definition}[MRAE-стійке AEAD]
\label{def:mrae}
AEAD-схема $(\Enc, \Dec)$ є \emph{стійкою до повторного використання
nonce} (\MRAE{}, \emph{nonce-misuse resistance})~\cite{rogaway2006},
якщо два повідомлення, зашифровані одним і тим самим ключем та
nonce, не дають опоненту жодної переваги, окрім тривіальної
відмінності фактів рівності/нерівності відкритих текстів.
Використовувана \Aura{} схема \AEAD{} (RFC~8452~\cite{rfc8452}) задовольняє
MRAE: випадкове або синтетичне (SIV) утворення внутрішнього лічильника
з ключа та AAD робить шифротекст ідентичним лише при ідентичних
відкритих текстах.
\end{definition}

\begin{remark}[Чому саме \AEAD{}]
Стандартний AES-GCM має катастрофічну деградацію безпеки при повторному
використанні nonce (втрату ключа автентифікації). У ратчет-протоколі
такий ризик не теоретичний: помилка серіалізації стану, відкат бази,
fork-процесу — все це може породити nonce-повтор. \AEAD{} (GCM-SIV) є
повноцінним AEAD у регулярному режимі та деградує до семантичної
безпеки лише за фактом повторного використання, без втрати
автентифікації.
\end{remark}

\subsection{Специфічна для \Aura{} нотація}\label{sec:aura-notation}

Уведемо такі позначення (\cref{tab:notation}):

\begin{table}[H]
\centering
\caption{Нотація специфічних для \Aura{} об'єктів та ключів.}\label{tab:notation}
\small
\begin{tabularx}{\textwidth}{@{}lX@{}}
\toprule
Символ & Опис \\
\midrule
$\IK_A^{x}, \IK_A^{ed}, \IK_A^{k}$ & довгострокові ключі сторони
  $A$: \DH{}-частина $\IK_A^{x}$ (X25519), Ed25519-ключ підпису
  $\IK_A^{ed}$ та \KEM{}-частина $\IK_A^{k}$ (Kyber-768 публ. ключ) \\
$\SPK_A$ & підписаний пре-ключ сторони $A$ (X25519) із підписом
  Ed25519 за допомогою $\IK_A^{ed}$ \\
$\OPK_A^{(i)}$ & $i$-й одноразовий пре-ключ сторони $A$ (X25519) \\
$\EK_A$ & ефемерний \DH{}-ключ ініціатора у X3DH \\
$\rk_n$ & кореневий ключ після $n$-го кроку ратчета \\
$\ck^{\rightarrow}, \ck^{\leftarrow}$ & ланцюгові ключі для напрямків
  «надсилання» та «прийом» \\
$\mk_{n,j}$ & ключ повідомлення з ланцюга епохи $n$, індекс $j$ \\
$\mdk_n$ & ключ шифрування метаданих епохи $n$ \\
$\epochv$ & номер ратчетної епохи (цілий, монотонно зростає) \\
$\mathit{pk\_new}_n, \mathit{sk\_new}_n$ & нові Kyber-ключі, згенеровані
  на $n$-му кроці ратчета \\
\bottomrule
\end{tabularx}
\end{table}

До тегових рядків належать, зокрема:
\begin{center}
\small
\path|Aura-X3DH|,\quad \path|Aura-Hybrid-X3DH|,\quad
\path|Aura-Hybrid-Ratchet|,\quad \path|Aura-MetadataKey|.
\end{center}
Рядки, використовувані як $\mathit{info}$-параметри HKDF, є
константами протоколу; повний перелік доменів подано в
\cref{tab:info-strings-ua}.

% ═════════════════════════════════════════════════════════════
% 4. АРХІТЕКТУРА ПРОТОКОЛУ
% ═════════════════════════════════════════════════════════════
\section{Архітектура протоколу}\label{sec:protocol}

\subsection{Огляд життєвого циклу сесії}\label{sec:overview}

Життєвий цикл двосторонньої сесії \Aura{} проходить чотири фази,
зображені у \cref{fig:lifecycle}:

\begin{enumerate}
  \item \textbf{Реєстрація.} Кожен користувач генерує довгостроковий
    \DH{}-ключ ідентифікації $\IK^x$ (X25519), Ed25519-ключ для підпису
    $\IK^{ed}$,
    \KEM{}-ключ ідентифікації $\IK^k$ (Kyber-768), підписаний пре-ключ
    $\SPK$ та пул одноразових пре-ключів $\{\OPK^{(i)}\}$. Усі публічні
    компоненти разом із підписами публікуються на сервер-каталог.
  \item \textbf{Рукостискання.} Ініціатор $A$ завантажує бандл $B$ із
    сервера та виконує гібридний X3DH (\S\ref{sec:handshake}). Результат —
    спільний кореневий ключ $\rk_0$, з якого деривуються початкові
    ланцюговий ключ $\ck_0$ та ключ метаданих $\mdk_0$.
  \item \textbf{Обмін повідомленнями.} Кожне повідомлення в одному
    напрямку обчислюється шляхом одностороннього ratchet ланцюгового
    ключа. На кожній зміні напрямку запускається гібридний Double
    Ratchet (\S\ref{sec:ratchet}): свіжий \DH{}- та \KEM{}-обмін.
  \item \textbf{Запечатане збереження стану.} Стан сесії періодично
    запечатується (AES-GCM-SIV з ключем, що зберігається у KeyStore
    платформи) і записується у локальну базу даних
    (\S\ref{sec:sealed}).
\end{enumerate}

\begin{figure}[H]
\centering
\begin{minipage}{0.95\textwidth}
\small
\begin{verbatim}
   Реєстрація           Рукостискання          Обмін повідомленнями
   -----------         ---------------        ----------------------
   IK_x, IK_k          [A] завант. бандла     ck_0 -> ck_1 -> ck_2 -> ...
   SPK, sig            [A->B] init+ct         (одно-напр. ratchet ck)
   {OPK^(i)}           [B] X3DH+Kyber->rk_0      |
        |              [B->A] ack                |  (зміна напрямку)
        +----> публ.   [A] фініш->rk_0           v
               на серв.                       Double Ratchet:
                                              X25519 свіжий + Kyber свіжий
                                              + augmented HKDF
                                              -> rk_{n+1}, ck_{n+1}, mdk_{n+1}
                                                 |
                                                 v
                                              Запечатаний стан
                                              (sealed; KEK у KeyStore)
\end{verbatim}
\end{minipage}
\caption{Життєвий цикл сесії \Aura{}.}\label{fig:lifecycle}
\end{figure}

\subsection{Пре-ключовий бандл}\label{sec:bundle}

Пре-ключовий бандл $B_A$ сторони $A$ — це опублікований у серверному каталозі
набір:
\begin{equation}\label{eq:prekey-bundle}
  B_A = \big(\IK_A^x, \, \IK_A^{ed}, \, \sigma_A, \, \IK_A^k, \, \SPK_A, \,
            \mathsf{Sign}_{\IK_A^{ed}}(\SPK_A), \,
            \{\OPK_A^{(i)}\}_{i \in I}\big),
\end{equation}
де $\sigma_A$ — Ed25519-підпис, який власна сторона $A$ ставить
над $\IK_A^x$ ключем $\IK_A^{ed}$ (доказ володіння та зв'язування
ідентичності); далі $\SPK_A$ підписується тим же Ed25519-ключем;
кожен
$\OPK_A^{(i)}$ — це нова X25519-пара з ідентифікатором $i$.

Сервер-каталог \emph{не} є довіреним: бандл валідується клієнтом
повністю, з перевіркою всіх підписів, ненульовості ключів та
відсутності точок малого порядку (RFC~7748~\S6.1). При запиті
бандла сервер атомарно виключає взятий $\OPK^{(i)}$ із пулу, аби
запобігти повторному використанню одноразового пре-ключа двома
ініціаторами.

\subsection{Гібридний X3DH}\label{sec:handshake}

Гібридний X3DH-обмін \Aura{} поєднує чотири DH-секрети
(класична X3DH-структура) із одним секретом інкапсуляції \KEM{}.

\paragraph{Дії ініціатора $A$.}
\begin{enumerate}
  \item Генерує ефемерний \DH{}-ключ: $(\EK_A^{\mathit{priv}}, \EK_A) \gets \DH.\KeyGen()$.
  \item Інкапсулює до \KEM{}-ключа ідентифікації $\IK_B^k$:
    $(c, K_{\KEM}) \gets \Encaps(\IK_B^k)$.
  \item Обчислює X3DH-секрети, де $\mathit{DH}_4$ використовується лише
    за наявності одноразового пре-ключа:
    \begin{align}
      \mathit{DH}_1 &= \DH(\IK_A^{x,\mathit{priv}}, \SPK_B), \label{eq:x3dh-dh1}\\
      \mathit{DH}_2 &= \DH(\EK_A^{\mathit{priv}}, \IK_B^x), \label{eq:x3dh-dh2}\\
      \mathit{DH}_3 &= \DH(\EK_A^{\mathit{priv}}, \SPK_B), \label{eq:x3dh-dh3}\\
      \mathit{DH}_4 &= \DH(\EK_A^{\mathit{priv}}, \OPK_B^{(i)})
        \quad\text{(якщо $\OPK_B^{(i)}$ наявний).} \label{eq:x3dh-dh4}
    \end{align}
  \item Формує класичний X3DH-вхід:
    \begin{equation}\label{eq:hybrid-ikm}
      \mathit{IKM} = \underbrace{\mathit{0xFF}^{32}}_{\text{padding}}
      \concat \mathit{DH}_1 \concat \mathit{DH}_2 \concat \mathit{DH}_3
      [\concat \mathit{DH}_4].
    \end{equation}
  \item Деривує класичний секрет і гібридний кореневий ключ:
    \begin{align}
      \mathit{cs}
      &\gets \HKDF(\mathit{IKM},\,32,\,\varepsilon,\,
        \texttt{Aura-X3DH}), \label{eq:x3dh-classical-secret}\\
      \rk_0
      &\gets \HKDF(K_{\KEM},\,32,\,\mathit{cs},\,
        \texttt{Aura-Hybrid-X3DH}). \label{eq:hybrid-derive}
    \end{align}
  \item Деривує ідентифікатор сесії, ключ метаданих і ключі підтвердження:
    \begin{align}
      \mathit{sid} &\gets \HKDF(\rk_0,\,16,\,c,\,
        \texttt{Aura-SessionId}),\\
      \mathit{ctx}_{md} &= \mathrm{sort}(\EK_A,\SPK_B)\concat \mathit{sid},\\
      \mdk_0 &\gets \HKDF(\rk_0,\,32,\,\mathit{ctx}_{md},\,
        \texttt{Aura-MetadataKey}),\\
      \mathit{kc}_I &\gets \HKDF(\rk_0,\,32,\,\varepsilon,\,
        \texttt{Aura-KeyConfirm-I}),\\
      \mathit{kc}_R &\gets \HKDF(\rk_0,\,32,\,\varepsilon,\,
        \texttt{Aura-KeyConfirm-R}).
    \end{align}
  \item Обчислює transcript-hash $\mathit{th}$ над канонічним бандлом і
    повідомленням ініціації з очищеним полем \texttt{key\_confirmation\_mac},
    надсилає $B$ ініціюючий конверт
    $(c, \EK_A, \IK_A^x, \IK_A^{ed}, \IK_A^k, i,
    \sigma_A^{\mathit{bind}}, \mathit{mac}_I)$, де
    $\sigma_A^{\mathit{bind}}$ підписує доменно відокремлене повідомлення
    \texttt{Aura-Handshake-Init-Identity-Binding} $\concat IK_A^{ed}\concat
    IK_A^x\concat IK_A^k$,
    $\mathit{mac}_I=\HMAC(\mathit{kc}_I,\mathit{th})$, і зберігає
    очікуваний $\mathit{mac}_R=\HMAC(\mathit{kc}_R,\mathit{th})$ для ACK.
\end{enumerate}

\paragraph{Дії респондента $B$.}
$B$ декапсулює $K_{\KEM} \gets \Decaps(\IK_B^{k,\mathit{priv}}, c)$,
обчислює ті самі X3DH-секрети у дзеркальному порядку та
застосовує \cref{eq:hybrid-ikm,eq:hybrid-derive}. Після перевірки
$\mathit{mac}_I$ респондент надсилає ACK із
$\mathit{mac}_R=\HMAC(\mathit{kc}_R,\mathit{th})$; ініціатор завершує
рукостискання лише після перевірки цього тегу. Сторони отримують один
і той же $\rk_0$.

\paragraph{Префікс $0\mathtt{xFF}^{32}$.}
Цей фіксований 32-байтовий префікс у $\mathit{IKM}$ забезпечує
доменну сепарацію між X3DH і будь-яким можливим іншим використанням
DH-секретів. Це стандартний підхід Signal X3DH; ми зберігаємо його
для зворотної сумісності аналізу.

\begin{theorem}[Конструкція IKM як гібридний комбінатор —
неформальне формулювання]
\label{thm:hybrid-ikm-informal}
Якщо класичний секрет $\mathit{cs}$ або KEM-секрет $K_{\KEM}$ є
непередбачуваним для опонента (\GapCDH{} для DH-частини або \INDCCA{} для
KEM-частини), то $\rk_0$
псевдовипадковий у моделі ROM з \dualPRF{} HKDF.
\end{theorem}

Формальне доведення наведено у \cref{thm:combiner}.

\subsection{Гібридний Double Ratchet}\label{sec:ratchet}

Основна відмінність гібридного Double Ratchet \Aura{} від класичного
Double Ratchet полягає в тому, що кожна зміна напрямку оновлює і
\DH{}-, і \KEM{}-матеріал із генерацією свіжої пари Kyber-ключів.

\subsubsection{Ініціалізація сесії після handshake}

Після того як гібридний X3DH дав $\rk_0$ і KEM-секрет, сесія виконує
початковий гібридний ratchet-крок і розділяє початкові ланцюги:
\begin{align}
  \mathit{dh}_{\mathit{init}}
    &= \DH(d_{\mathit{local}}, D_{\mathit{remote}}), \label{eq:init-dh-ua}\\
  \mathit{hybrid\_ikm}
    &= \mathit{dh}_{\mathit{init}} \concat K_{\KEM}, \label{eq:init-hybrid-ikm-ua}\\
  \rk_1
    &= \HKDF(\mathit{hybrid\_ikm},\,32,\,\rk_0,\,
      \texttt{Aura-Hybrid-Ratchet})[0{:}32], \label{eq:init-root-key-ua}\\
  (\ck_{\mathit{send}},\,\ck_{\mathit{recv}})
    &= \HKDF(\rk_1,\,64,\,\varepsilon,\,
      \texttt{Aura-ChainInit}). \label{eq:init-chain-keys-ua}
\end{align}
Ініціатор використовує перші 32~байти як $\ck_{\mathit{send}}$, а
останні 32~байти як $\ck_{\mathit{recv}}$; респондент міняє ці ролі
місцями.

\subsubsection{Опис кроку гібридного ратчета}\label{sec:ratchet-step}

Розглянемо момент, коли сторона $A$, яка раніше була у режимі
«прийому» (\emph{receiving}), має надіслати перше повідомлення нового
напрямку. $A$ виконує крок ратчета:

\begin{algorithm}[H]
\caption{Гібридний крок ратчета на стороні $A$ (відправка).}\label{alg:ratchet-step}
\KwIn{Поточний стан $\sigma = (\rk_n, \mathit{dh\_remote}, \mathit{pk\_remote})$}
\KwOut{Оновлений стан $\sigma' = (\rk_{n+1}, \ck_{n+1}^{\rightarrow},
\mdk_{n+1}, \mathit{dh\_local}_{n+1}, \mathit{pk\_local}_{n+1},
\mathit{sk\_local}_{n+1})$,
заголовок $h$}

$(\mathit{dh\_priv}_{n+1}, \mathit{dh\_pub}_{n+1}) \gets \DH.\KeyGen()$\;
$\mathit{dh\_ss} \gets \DH(\mathit{dh\_priv}_{n+1}, \mathit{dh\_remote})$\;
$(c_{n+1}, k_{\KEM}) \gets \Encaps(\mathit{pk\_remote})$
  \Comment*{використовуємо поточний публ. KEM-ключ контрагента}
$(\mathit{pk\_local}_{n+1}, \mathit{sk\_local}_{n+1}) \gets \KEM.\KeyGen()$
  \Comment*{свіжа пара KEM-ключів для наступного раунду}
$\mathit{ikm} \gets \mathit{dh\_ss} \concat k_{\KEM}$\;
$\mathit{info} \gets \texttt{Aura-Hybrid-Ratchet} \concat \mathit{pk\_local}_{n+1}$
  \Comment*{доповнений HKDF: новий KEM-ключ у info}
$(\rk_{n+1} \concat \ck_{n+1}^{\rightarrow} \concat \mdk_{n+1})
  \gets \HKDF(\mathit{ikm}; \mathit{salt} = \rk_n,\, \mathit{info})$\;
$h \gets (\mathit{dh\_pub}_{n+1}, \mathit{pk\_local}_{n+1}, c_{n+1},
  \epochv_{n+1}, \mathit{prev\_chain\_len})$\;
\Return{$\sigma' , h$}
\end{algorithm}

\paragraph{Дзеркальні дії на стороні $B$.}
Сторона $B$, отримавши заголовок $h$, декапсулює
$k_{\KEM} \gets \Decaps(\mathit{sk}_{B}^{\mathit{current}}, c_{n+1})$
власним поточним KEM-секретом,
обчислює $\mathit{dh\_ss}$ симетрично, формує те саме $\mathit{info}$
(із прийнятим $\mathit{pk\_local}_{n+1}$) та деривує той самий $\rk_{n+1}$.

\subsubsection{Доповнений HKDF: зв'язування KEM-ключа}\label{sec:augmented}

У рядку $\mathit{info} := \texttt{Aura-Hybrid-Ratchet}
\concat \mathit{pk\_local}_{n+1}$ у \cref{alg:ratchet-step} до
доменного тегу додається новий KEM-публічний ключ. Далі цю техніку
називатимемо \emph{augmented HKDF binding}.

\begin{remark}[Сенс доповненого HKDF]
HKDF — це \dualPRF{}: при фіксованих $\mathit{ikm}, \mathit{salt}$ вихід
залежить детерміновано від $\mathit{info}$. Якщо опонент-у-середині
підмінює $\mathit{pk\_local}_{n+1}$ на свій варіант $\mathit{pk}'$
у заголовку $h$, тоді у $A$ і $B$ обчислиться різний
$\mathit{info}$ і, як наслідок, різний $\rk_{n+1}$. Перше ж AEAD-розшифрування
наступного повідомлення (з тегом цілісності) провалиться, і атака
виявляється негайно --- без додаткового round-trip і без цифрового
підпису над KEM-ключем.
\end{remark}

\begin{lemma}[Augmented HKDF — неформально]\label{lem:augmented-informal}
У моделі ROM, при припущенні, що HKDF поводиться як випадковий оракул,
для будь-якого \ppt{} опонента $\advA$ ймовірність того, що
обчислений на стороні $B$ ключ $\rk_{n+1}'$ збігається з ключем
$\rk_{n+1}$, обчисленим стороною $A$, за умови, що
$\mathit{pk}' \neq \mathit{pk\_local}_{n+1}$, є знехтуваною
($2^{-256}$ для SHA-256-based HKDF).
\end{lemma}

Аргумент подано у \cref{lem:augmented-informal}; формальна версія
входить до доведення AKE.

\paragraph{Роль зв'язування KEM-ключа.}
Класичний Double Ratchet домагається неявної автентифікації лише
\DH{}-ключа через те, що той бере участь у деривації. Для KEM-ключа
ця властивість не виникає сама собою, бо в стандартному гібридному
комбінаторі (\cref{eq:hybrid-ikm}) KEM-ключ як \emph{публічний} об'єкт
не з'являється — з'являється лише похідний $k_{\KEM}$. Доповнений
HKDF явно включає публічний KEM-ключ до входу деривації,
що рівнозначно вимозі узгодженості заголовка з тілом.

\subsection{Дворівневе AEAD з шифруванням метаданих}\label{sec:two-layer}

Дворівнева конструкція AEAD \Aura{} відокремлює шифрування метаданих
від шифрування корисного навантаження.

Кожне AAD-поле містить ідентифікаційне зв'язування сесії:
\begin{equation}\label{eq:identity-binding}
  \mathit{ibh} =
  \SHA(\texttt{Aura-Identity-Binding}
  \concat \mathrm{sort}(\mathit{ed}_I,\mathit{ed}_R)
  \concat \mathrm{sort}(\mathit{x}_I,\mathit{x}_R)).
\end{equation}
Ed25519- та X25519-публічні ключі сортуються лексикографічно перед
конкатенацією, щоб одна й та сама пара сторін мала однаковий binding
незалежно від ролі ініціатора.

\paragraph{Рівень 1 — метадані конверта.}
Заголовок конверта (тип повідомлення, лічильник $j$,
ідентифікатор кореляції) серіалізується як
бінарний об'єкт $\mathit{meta}$ та шифрується незалежним ключем
$\mdk_n$ (який деривується одночасно з $\rk_n$ та $\ck_n$
з виходу HKDF на кроці ратчета, див. \cref{alg:ratchet-step}):
\begin{equation}\label{eq:metadata-aead}
  \mathit{ct\_meta} \gets \Enc_{\mdk_n}(N_{\mathit{meta}},\,
  \mathit{meta},\, \mathit{AAD}_{\mathit{meta}}).
\end{equation}
Тут
\begin{equation}\label{eq:metadata-aad}
  \mathit{AAD}_{\mathit{meta}}
  = \mathit{sid} \concat \mathit{ibh} \concat \epochv_n \concat v,
\end{equation}
де $v$ — версія протоколу.

\paragraph{Рівень 2 — корисне навантаження.}
Із ланцюгового ключа $\ck_n$ деривуються ключ повідомлення та наступний
стан ланцюга:
\begin{equation}\label{eq:message-key-derive}
  \mk_{n,j} \gets \HKDF(\ck_{n,j},\,32,\,\varepsilon,\,
    \texttt{Aura-Msg}),
\end{equation}
\begin{equation}\label{eq:chain-key-derive}
  \ck_{n,j+1} \gets \HKDF(\ck_{n,j},\,32,\,\varepsilon,\,
    \texttt{Aura-Chain}).
\end{equation}
де $\ck_{n,j}$ — стан ланцюга після $j$-ї ітерації. Корисне навантаження
шифрується схемою \AEAD{}:
\begin{equation}\label{eq:payload-aead}
  \mathit{ct\_payload} \gets \Enc_{\mk_{n,j}}(N_{\mathit{payload}},\,
  \mathit{plaintext},\, \mathit{AAD}_{\mathit{payload}}),
\end{equation}
де
\begin{equation}\label{eq:payload-aad}
  \mathit{AAD}_{\mathit{payload}}
  = \mathit{sid} \concat \mathit{ibh} \concat \epochv_n \concat j \concat v.
\end{equation}
Отже, обидва AEAD-рівні прив'язані до сесії, ролей сторін, епохи та
версії протоколу; payload-рівень додатково прив'язано до індексу
повідомлення.

Конверт у мережі несе відкрито лише службові поля, потрібні для маршрутизації
і дешифрування:
\begin{itemize}
  \item 12-байтовий $\mathit{header\_nonce}$ для розшифрування metadata-рівня;
  \item шифротекст метаданих $\mathit{ct\_meta}$;
  \item шифротекст корисного навантаження $\mathit{ct\_payload}$;
  \item номер епохи $\epochv_n$;
  \item на ratchet-кроці — $\mathit{dh\_pub}_n$, $c_n$,
    $\mathit{pk\_local}_n$ та довжину попереднього ланцюга.
\end{itemize}
Лічильник повідомлення, payload nonce, тип envelope і correlation id
залишаються всередині зашифрованих метаданих; відкритий header nonce не
кодує порядку повідомлень.

\paragraph{Ротація ключа метаданих.}
Ключ $\mdk_n$ повністю змінюється на кожній зміні напрямку ратчета
(епоха $n \to n+1$). У межах однієї епохи ключ $\mdk_n$ використовується
для всіх повідомлень, але nonce-простір AEAD надійно сепарується
(\S\ref{sec:nonce}), а MRAE-стійкість \AEAD{} забезпечує безпеку
навіть у разі катастрофічної помилки nonce-керування.

\paragraph{Кешування $\mdk$ для out-of-order доставки.}
Реалізація підтримує впорядкований кеш до 100~останніх ключів
метаданих (константа $\textsf{MAX\_CACHED\_METADATA\_KEYS} = 100$),
обмежений за віком 32~епохами
($\textsf{MAX\_METADATA\_KEY\_EPOCH\_AGE} = 32$). Це дозволяє
розшифрувати метадані повідомлень, що прийшли поза порядком, не
поступаючись вікном прямої секретності для метаданих.

\paragraph{Чому AEAD саме у GCM-SIV-режимі.}
У ратчет-протоколі катастрофічна помилка nonce-керування цілком
реалістична: повторний крок ратчета після виключення процесу до
збереження стану, відновлення з резервної копії бази, мережна
дублікація. AEAD-GCM поза SIV-режимом за повторного nonce втрачає
автентифікаційний ключ і дозволяє підробку повідомлень. GCM-SIV
у такому ж сценарії зберігає misuse-resistant межі конфіденційності й
цілісності, але повтор nonce може розкрити рівність однакових пар
AAD/plaintext. Ця модель загроз мотивує використання GCM-SIV замість
звичайного GCM у протокольному envelope.

\subsection{Захист від повторів та керування nonce}\label{sec:replay}\label{sec:nonce}

\paragraph{Структура nonce.}
12-байтовий nonce AEAD конструктивно складається з:
\begin{equation}\label{eq:nonce-layout}
  N = \underbrace{\mathit{prefix}_{8B}}_{\text{CSPRNG}}
      \concat \underbrace{\mathit{counter}_{2B}}_{\text{лічильник чанку}}
      \concat \underbrace{\mathit{index}_{2B}}_{j \bmod 2^{16}}.
\end{equation}
8-байтовий prefix генерується CSPRNG під час створення сесії або
ratchet-reset і зберігається у стані; 2-байтовий counter та 2-байтовий
index обмежені відповідно $2^{16}-1$. Реалізація форсує новий
ratchet-крок при вичерпанні ланцюга повідомлень ($N_{\max}=10{,}000$)
і викликає callback попередження, коли залишок nonce-бюджету падає до
10\% або нижче. Metadata-рівень використовує незалежний 12-байтовий
\texttt{header\_nonce}.

\paragraph{Дедуплікація nonce.}
На приймальному боці підтримується bounded кеш до 2048~останніх побачених
$\langle$епоха, nonce$\rangle$-пар; повтор виявляється негайно.
Це поверх захисту, який дає AEAD-MRAE.

\paragraph{Дедуплікація ініціаційних повідомлень.}
Респондент додатково підтримує LRU-кеш до 16{,}384 хешів отриманих
\textsf{HandshakeInit} повідомлень
($\textsf{MAX\_SEEN\_HANDSHAKE\_INITS}$), щоб запобігти повторам
рукостискань.

\subsection{Запечатане збереження стану}\label{sec:sealed}

Між викликами користувача протокол повинен серіалізувати свій стан
до локального сховища. \Aura{} застосовує конструкцію
\emph{запечатаного} (sealed) стану з двошаровим шифруванням:

\begin{enumerate}
  \item Генерується ефемерний 32-байтовий ключ шифрування даних
    $\mathit{DEK} \getsr \{0,1\}^{256}$.
  \item Стан серіалізується в protobuf $\sigma_{\mathit{pb}}$ і
    шифрується: $\mathit{ct}_{\sigma} \gets \Enc_{\mathit{DEK}}(N_{\sigma},
    \sigma_{\mathit{pb}}, \mathit{AAD}_{\sigma})$.
  \item $\mathit{DEK}$ шифрується ключем шифрування ключів
    $\mathit{KEK}$, що зберігається у платформному
    KeyStore/Keychain/TPM: $\mathit{ct}_{\mathit{DEK}} \gets
    \Enc_{\mathit{KEK}}(N_{\mathit{DEK}}, \mathit{DEK}, \mathit{AAD}_{\mathit{DEK}})$.
  \item На диск пишеться кортеж $(v, \mathit{external\_counter}, N_{\mathit{DEK}},
    \mathit{ct}_{\mathit{DEK}}, N_{\sigma}, \mathit{ct}_{\sigma})$.
\end{enumerate}

\paragraph{Зовнішній лічильник проти відкату.}
Поле $\mathit{external\_counter}$ — це монотонно зростаючий
номер запечатаних записів, що зберігається в окремому,
довіреному до запису сховищі (наприклад, секурном-енклаві,
WAL-сегменті бази даних, або у системному заголовку файла з
правами лише адміністратора). При відкритті стану викликач передає
\emph{мінімальний прийнятний} лічильник $\mathit{min\_external\_counter}$;
протокол відкидає стан, якщо $\mathit{external\_counter} <
\mathit{min\_external\_counter}$.

\begin{remark}[Чому строге порівняння]\label{rem:strict}
Раніша версія перевіряла $\mathit{counter} \le \mathit{min}$ як
ознаку відкату, проте це породжувало хибне виявлення на холодному
старті iOS-клієнта: повторне відкриття месенджера без створення
жодного нового стану лишало $\mathit{counter} == \mathit{min}$ і
спричиняло помилку $\mathsf{AuraErrorInvalidState}$
(FFI код~16). Поточна версія використовує строге $<$, що:
\emph{(а)}~допускає відновлення з ідентичного збереженого стану як
законний випадок; \emph{(б)}~продовжує блокувати справжній
відкат, де $\mathit{counter}$ менший за мінімальний. Доповнено
ідемпотентністю виклику $\textsf{note\_successful\_restore}$ при
рівному лічильнику, та зміною оператора
$\textsf{min\_import\_counter}() = \max(\textsf{max\_restored},
\textsf{latest\_issued})$ для коректного виявлення reverse-downgrade.
\end{remark}

\paragraph{Цілісність стану.}
Поверх AEAD-цілісності, по всьому збереженому об'єкту обчислюється
HMAC-SHA256 з ключем, деривованим із поточного root key:
$\HKDF(\rk;\mathit{info}=\texttt{Aura-State-HMAC})$. Це дає додатковий
рівень захисту проти tampering із серіалізованим станом поза
AES-GCM-SIV. Зовнішній monotonic counter sealed state лишається окремим
механізмом anti-rollback.

\subsection{Архітектурні інваріанти реалізації}\label{sec:impl-invariants}

Нижче наведено реалізаційні інваріанти, які використовує proof-to-code
map. Доведення оперує абстрактними оракулами та ключами, тоді як
реалізація має обмежені буфери, конкретні protobuf-поля, кеші, коди
помилок та FFI-межі. Порушення цих інваріантів змінювало б не лише
реалізацію, а й модель безпеки.

\paragraph{Доменна сепарація.}
Кожна семантична роль деривації має окрему мітку; повний перелік подано
в \cref{tab:info-strings-ua}. Цей набір не є косметичним: повторне
використання однієї мітки для різних ключів зливало б простори
повідомлень HKDF і робило б частину доведень некоректною. Наприклад,
ключ повідомлення і ключ метаданих виводяться з одного ратчетного
виходу, але різними доменами; саме це дозволяє трактувати їх як
незалежні робочі ключі в теоремі про конфіденційність та цілісність.

\begin{table}[H]
\centering
\caption{HKDF/domain-separation рядки, використані у \Aura{}.}
\label{tab:info-strings-ua}
\small
\begin{tabular}{@{}ll@{}}
\toprule
Рядок & Контекст \\
\midrule
\texttt{Aura-X3DH} & класична деривація X3DH \\
\texttt{Aura-Hybrid-X3DH} & гібридний combiner у handshake \\
\texttt{Aura-Hybrid-Ratchet} & гібридний ratchet-крок, augmented із $kpk$ \\
\texttt{Aura-DH-Ratchet} & класичний fallback DH-ратчет \\
\texttt{Aura-ChainInit} & початковий поділ send/recv chain keys \\
\texttt{Aura-Chain} & перехід symmetric chain key \\
\texttt{Aura-Msg} & деривація message key \\
\texttt{Aura-SessionId} & ідентифікатор сесії \\
\texttt{Aura-MetadataKey} & ключ шифрування метаданих \\
\texttt{Aura-State-HMAC} & HMAC-ключ sealed-state integrity \\
\texttt{Aura-KeyConfirm-I} & key confirmation ініціатора \\
\texttt{Aura-KeyConfirm-R} & key confirmation респондента \\
\texttt{Aura-Handshake-Transcript} & transcript hash \\
\texttt{Aura-Identity-Binding} & identity binding hash \\
\texttt{Aura-PQ-Hybrid::} & префікс солі гібридного KEM-combiner \\
\bottomrule
\end{tabular}
\end{table}

\paragraph{Фіксовані розміри як частина протоколу.}
Реалізація фіксує 32-байтові кореневі, ланцюгові, повідомленнєві та
metadata-ключі, 12-байтовий AEAD nonce, 32-байтові X25519 публічні
ключі, а також розміри ML-KEM-768: 1184~байти публічного ключа,
1088~байтів ciphertext і 2400~байтів секретного ключа. Ці числа
входять до інваріантів парсингу: повідомлення з іншими довжинами
має відхилятися до криптографічної операції. Інакше реалізація
переносила б помилки формату в криптографічний шар, де вони важче
розрізняються і можуть перетворитись на oracle-поведінку.

\paragraph{Обмежені кеші.}
Кеш пропущених message keys, кеш metadata-ключів, LRU-кеш побачених
nonce та кеш ініціацій рукостискання мають явні верхні межі. Це
одночасно інженерна та криптографічна вимога. Інженерно вона
запобігає DoS через нескінченне накопичення out-of-order стану.
Криптографічно вона робить модель ближчою до реального опонента:
повтор або дуже старе повідомлення не «чекає» в системі безмежно,
а або обробляється в допустимому вікні, або відкидається як застаріле.
Саме тому вікна \textsf{MAX\_SEEN\_NONCES},
\textsf{MAX\_SEEN\_HANDSHAKE\_INITS} та
\textsf{MAX\_CACHED\_METADATA\_KEYS} є параметрами протоколу, а не
випадковими константами реалізації.

\paragraph{Атомарність кроку ратчета.}
Крок ратчета змінює одразу кілька величин: кореневий ключ, ланцюговий
ключ напряму, ключ метаданих, локальну KEM-пару та заголовок наступної
епохи. Реалізаційний інваріант полягає в тому, що ці зміни мають
фіксуватись як одна транзакція: часткове оновлення, після якого стан
збережено, але заголовок або metadata-key ще не синхронізовано,
створює клас помилок, який формальна модель не покриває. Тому sealed
state і rollback-лічильник у \cref{sec:sealed} розглядаються як
частина протоколу, а не як зовнішня зручність клієнта.

\paragraph{Межа FFI.}
FFI-шар не повинен змінювати криптографічну семантику: він лише
перетворює байтові буфери та коди помилок між Rust і зовнішнім
середовищем. Усі перевірки довжин, стану сесії, replay-вікон і
rollback-лічильників мають відбуватися до повернення успішного коду
через FFI. Така вимога особливо важлива для мобільних клієнтів, де
відновлення зі сховища, фонове завершення процесу та повторний запуск
часто трапляються між двома повідомленнями.

% ═════════════════════════════════════════════════════════════
% 1. МОДЕЛЬ БЕЗПЕКИ
% ═════════════════════════════════════════════════════════════
\section{Формальна модель безпеки}\label{sec:model}

Адаптуємо модель Канетті--Кравчика~\cite{canetti2001} та її
розширену \eCK{}-версію~\cite{lamacchia2007} до гібридного
постквантового сценарію.
Модель \eCK{} обрана через те, що вона надає опоненту широкий
практично-релевантний набір можливостей: окремий доступ до
довгострокових ключів та до ефемерного стану сесії, у будь-якій
комбінації, окрім тривіально-виграшних.

\subsection{Сторони та сесії}\label{sec:parties}

Розглядаємо щонайбільше $N$ сторін $P_1, \ldots, P_N$. Кожна
сторона $P_i$ має:
\begin{itemize}
  \item довгостроковий ідентифікаційний \DH{}-ключ $\IK_i$ (X25519);
  \item довгостроковий \KEM{}-ключ $(\mathit{kpk}_i, \mathit{ksk}_i)$
    (Kyber-768);
  \item ротаційні підписані пре-ключі $\SPK_i$ (X25519 + Ed25519-підпис).
\end{itemize}
\emph{Сесія} $\pi_i^s$ — це $s$-й екземпляр протоколу, який виконує
сторона $P_i$ з наміченим партнером $P_j$.

\subsection{Оракули опонента}\label{sec:oracles}

\begin{definition}[Оракули гри AKE-безпеки]\label{def:oracles}
Опонент $\advA$ взаємодіє з протоколом через оракули:
\begin{itemize}
  \item $\mathsf{Send}(\pi_i^s, m)$ — доставляє повідомлення $m$
    сесії $\pi_i^s$; повертає відповідь протоколу. Моделює повний
    контроль над мережею (Dolev--Yao).
  \item $\mathsf{Corrupt}(P_i)$ — повертає всі довгострокові секретні
    ключі $P_i$ (X25519-секрет $\mathit{ik}_i$, Kyber-секрет
    $\mathit{ksk}_i$). Моделює повну компрометацію довгострокового
    ключового матеріалу.
  \item $\mathsf{RevealState}(\pi_i^s)$ — повертає повний стан сесії
    $\pi_i^s$ (кореневий ключ, ланцюгові ключі, ключі метаданих,
    приватний \DH{}-ключ ратчета, приватний Kyber-ключ ратчета, стан
    nonce). Моделює компрометацію стану в конкретний момент часу.
  \item $\mathsf{RevealSessionKey}(\pi_i^s)$ — повертає ключ сесії
    (кореневий ключ $\rk_e$ поточної епохи).
  \item $\mathsf{Test}(\pi_i^s)$ — якщо біт виклику $b = 0$, повертає
    справжній ключ сесії; якщо $b = 1$ — рівномірно випадковий ключ
    тієї ж довжини. Дозволено лише один запит $\mathsf{Test}$.
\end{itemize}
\end{definition}

\subsection{Партнерство}\label{sec:partnering}

Дві сесії $\pi_i^s$ і $\pi_j^t$ є \emph{партнерами}, якщо вони
поділяють той самий хеш транскрипту $\mathit{th}$, обчислений як
HMAC-SHA256 над довжино-префіксованою конкатенацією пре-ключового
бандлу та повідомлення ініціації рукостискання.

\subsection{Свіжість}\label{sec:freshness}

\begin{definition}[Свіжість]\label{def:freshness}
Сесія $\pi_i^s$ з партнером $\pi_j^t$ є \emph{свіжою}, якщо не
виконується жодна з умов:
\begin{enumerate}
  \item $\advA$ викликав $\mathsf{RevealSessionKey}$ або
    $\mathsf{RevealState}$ для $\pi_i^s$ чи $\pi_j^t$ до запиту
    $\mathsf{Test}$;
  \item $\advA$ викликав і $\mathsf{Corrupt}(P_i)$, і
    $\mathsf{RevealState}(\pi_i^s)$ (тобто має водночас довгостроковий
    ключ та ефемерний стан однієї сторони — тривіальний виграш);
  \item $\advA$ викликав $\mathsf{Corrupt}(P_j)$ \emph{до} створення
    сесії $\pi_j^t$ \textbf{і} жоден одноразовий пре-ключ не було
    використано (активна імперсонація без OPK-захисту).
\end{enumerate}
\end{definition}

\subsection{Означення стійкості}\label{sec:defs}

\begin{definition}[\eCK{}-AKE-стійкість]\label{def:ake}
Протокол $\Pi$ є \eCK{}-AKE-стійким, якщо для всіх \ppt{} опонентів
$\advA$:
\begin{equation}\label{eq:eck-advantage}
  \adv^{\eCK}_{\Pi}(\advA) =
  \Big|\Pr[\advA \text{ виграє}] - \tfrac{1}{2}\Big|
  \le \negl(\secpar),
\end{equation}
де $\advA$ виграє, якщо правильно вгадує біт $b$ для свіжої
тестованої сесії.
\end{definition}

\begin{definition}[Пряма секретність]\label{def:fs}
Протокол $\Pi$ забезпечує пряму секретність, якщо компрометація
довгострокових ключів $(\IK_X, \mathit{ksk}_X)$ у момент~$t$ не
дозволяє розшифрувати повідомлення, зашифровані до~$t$ у сесіях,
що завершили обмін ключами до компрометації.
\end{definition}

\begin{definition}[$k$-крокова постскомпрометаційна
безпека]\label{def:pcs}
Протокол $\Pi$ забезпечує $k$-крокову постскомпрометаційну безпеку,
якщо після повної компрометації стану на епосі~$e$ (через
$\mathsf{RevealState}$) повідомлення епох $\ge e + k$ є безпечними
за умови, що опонент не компрометує \emph{нові} ключі, згенеровані
після епохи~$e$.
\end{definition}

% ═════════════════════════════════════════════════════════════
% 2. ТЕОРЕМИ БЕЗПЕКИ
% ═════════════════════════════════════════════════════════════
\section{Теореми безпеки}\label{sec:theorems}

Усі шість теорем доводяться методом послідовності ігор (game-hopping)
із побудовою явних редукцій. Записуємо доведення у формі ескізів;
повні розгорнуті доведення доступні у супровідному технічному звіті
(\texttt{security-proof.tex})~\cite{aura-security-proof}.

\subsection{Теорема 1: псевдовипадковість гібридного кореневого ключа}\label{sec:thm1}

\begin{theorem}[Псевдовипадковість гібридного кореневого ключа]\label{thm:combiner}
Якщо \KEM{} є \INDCCA{}-стійким \textbf{або} \DH{} задовольняє
\GapCDH{} у моделі \ROM{}, то гібридний комбінатор \Aura{}
\begin{equation}\label{eq:combiner-root}
  \rk_0 = \HKDF\big(\mathit{ikm} = \mathit{ss}_{\KEM};\;
    \mathit{salt} = \mathit{cs};\;
    \mathit{info} = \texttt{Aura-Hybrid-X3DH}\big)
\end{equation}
видає псевдовипадковий кореневий ключ. Конкретно, для будь-якого
\ppt{} опонента $\advA$:
\begin{equation}\label{eq:combiner-bound}
  \adv^{\mathsf{rk}}_{\mathsf{Hybrid}}(\advA)
  \le
  \min\!\left\{
    \adv^{\INDCCA}_{\KEM}(\advB_1),
    4 \cdot \adv^{\GapCDH}_{\DH}(\advB_2)
  \right\}
  + \adv^{\dualPRF}_{\HKDF}(\advB_3)
  + \frac{q_H}{2^{256}},
\end{equation}
де $q_H$ — кількість запитів до випадкового оракула, а множник~4
відповідає чотирьом DH-обчисленням класичного X3DH-компонента.
\end{theorem}

\begin{proof}[Ескіз доведення]
Є два шляхи редукції. У KEM-шляху замінюємо Kyber-секрет
$\mathit{ss}_{\KEM}$ на випадковий; різниця обмежена
$\adv^{\INDCCA}_{\KEM}(\advB_1)$ через редукцію $\advB_1$, що вкладає
\INDCCA{}-челендж як $\mathit{kpk}_R$. У класичному шляху замінюємо
секрет $\mathit{cs}$ на випадковий; різниця обмежена
$4\cdot\adv^{\GapCDH}_{\DH}(\advB_2)$, бо редукція вгадує один із
чотирьох DH-компонентів X3DH. У кожному з двох шляхів \emph{хоча б
один} вхід HKDF є випадковим, тому властивість \dualPRF{} робить
$\rk_0$ псевдовипадковим. Беремо кращий із двох шляхів, що дає
\cref{eq:combiner-bound}. \emph{Ключове спостереження}: гібридний
комбінатор має гарантію «\textbf{або-або}» — він безпечний, якщо
безпечний \emph{принаймні один} із двох примітивів, тож компрометація
лише класичного DH (наприклад, через CRQC) не руйнує $\rk_0$.
\end{proof}

\subsection{Теорема 2: \eCK{}-AKE-стійкість}\label{sec:thm2}

\begin{theorem}[\eCK{}-AKE-стійкість]\label{thm:ake}
Протокол гібридного X3DH \Aura{} є \eCK{}-AKE-стійким. Для будь-якого
\ppt{} опонента $\advA$, що взаємодіє з не більш ніж $q_s$ сесіями:
\begin{align}
  \adv^{\eCK}_{\Aura}(\advA)
  &\le
  q_s \cdot \big(
    4 \cdot \adv^{\GapCDH}_{\DH}(\advB_1)
    + \adv^{\INDCCA}_{\KEM}(\advB_2) \notag\\
  &\qquad
    + \adv^{\dualPRF}_{\HKDF}(\advB_3)
    + 2 \cdot \adv^{\PRF}_{\HMAC}(\advB_4)\big)
    + \frac{q_s^2}{2^{128}} + \frac{q_H}{2^{256}} .
  \label{eq:ake-bound}
\end{align}
\end{theorem}

\begin{proof}[Ескіз доведення]
Доведення йде послідовністю ігор. Спочатку редукція вгадує тестовану
сесію серед $q_s$ (множник $q_s$). Далі чотири DH-внески X3DH
замінюються випадковими величинами з втратою
$4\cdot\adv^{\GapCDH}_{\DH}$; KEM-секрет замінюється через
\INDCCA{}-стійкість \KEM{}; вихід HKDF замінюється через
\dualPRF{}-властивість; два теги key confirmation
(\texttt{Aura-KeyConfirm-I/R}) замінюються через \PRF{}-стійкість
\HMAC{}. Доданок $q_s^2/2^{128}$ обмежує колізії 128-бітних
ідентифікаторів сесії, а $q_H/2^{256}$ — колізії випадкового
оракула.
\end{proof}

\subsection{Теорема 3: пряма секретність}\label{sec:thm3}

\begin{theorem}[Пряма секретність]\label{thm:fs}
За припущення \GapCDH{} для стертих ефемерних DH-ключів протокол
\Aura{} забезпечує класичну пряму секретність (\cref{def:fs}) навіть
після пізнішого розкриття довгострокових KEM-секретів. Початковий
KEM-компонент дає захист від harvest-now-decrypt-later лише доти, доки
відповідний KEM-секрет не розкрито.
\begin{equation}\label{eq:fs-bound}
  \adv^{\mathsf{FS}}_{\Aura}(\advA)
  \le
  \adv^{\GapCDH}_{\DH}(\advB_1)
  + \adv^{\dualPRF}_{\HKDF}(\advB_2).
\end{equation}
\end{theorem}

\begin{proof}[Ескіз доведення]
Кожен ключ повідомлення $\mk_{n,j}$ і наступний ланцюговий ключ
$\ck_{n,j+1}$ деривуються з $\ck_{n,j}$ через HKDF з різними
мітками \texttt{Aura-Msg} і \texttt{Aura-Chain}
(\cref{eq:message-key-derive,eq:chain-key-derive}). Через
\PRF{}-стійкість HKDF знання $\ck_{n,j+1}$ не дає $\ck_{n,j}$
(і тим паче $\mk_{n,j}$). Кожен крок ратчета перезаписує кореневий ключ через
$\rk_{n+1} = \HKDF(\cdot; \mathit{salt} = \rk_n)$, причому старий
$\rk_n$ та проміжний ключовий матеріал негайно затираються
(\texttt{zeroize}). Отже, компрометація довгострокових або
поточних ключів у момент~$t$ не відновлює минулих $\mk$.
Початковий KEM-ciphertext не рахується в межу FS після розкриття
$\mathit{ksk}$, бо записаний ciphertext можна декапсулювати. Його
квантова роль у цій статті — HNDL-захист без пізнішого розкриття
KEM-секрету та PCS після появи свіжих KEM-ключів у ратчеті.
\end{proof}

\subsection{Теорема 4: постскомпрометаційна безпека та її
асиметрія}\label{sec:thm4}

Наступна теорема формалізує результат~\textbf{Н4}.

\begin{theorem}[Асиметрична PCS]\label{thm:pcs}
Нехай опонент виконав $\mathsf{RevealState}(\pi_i^e)$ на епосі~$e$
(повна компрометація стану) і надалі не компрометує нові ключі. Для
гібридної частини беремо консервативну модель компрометації обох
кінців, де опонент також має KEM-секрет того peer-ключа, до якого
піде перша післякомпрометаційна інкапсуляція. Тоді протокол \Aura{}
забезпечує:
\begin{enumerate}
  \item[\textbf{(а)}] \textbf{1-крокову класичну PCS}: за припущенням
    \GapCDH{} повідомлення епохи $\ge e+1$ є безпечними:
    \begin{equation}\label{eq:pcs-classical-bound}
      \adv^{\mathsf{PCS\text{-}1}}_{\Aura}(\advA)
      \le
      \adv^{\GapCDH}_{\DH}(\advB_1)
      + \adv^{\dualPRF}_{\HKDF}(\advB_2);
    \end{equation}
  \item[\textbf{(б)}] \textbf{2-крокову гібридну PCS}: за припущенням
    \GapCDH{} та \INDCCA{} для \KEM{} повідомлення епохи $\ge e+2$
    є безпечними:
    \begin{equation}\label{eq:pcs-hybrid-bound}
      \adv^{\mathsf{PCS\text{-}2}}_{\Aura}(\advA)
      \le
      \adv^{\GapCDH}_{\DH}(\advB_1)
      + \adv^{\INDCCA}_{\KEM}(\advB_3)
      + 2 \cdot \adv^{\dualPRF}_{\HKDF}(\advB_2).
    \end{equation}
\end{enumerate}
Обидві редукції не мають множника вгадування сесії.
\end{theorem}

\begin{proof}[Ескіз доведення]
На епосі $e+1$ сторона генерує свіжу \DH{}-пару
(\cref{alg:ratchet-step}), тож класичний \DH{}-внесок
до $\rk_{e+1}$ є свіжим і
невідомим опоненту — звідси~\textbf{(а)}. Однак на епосі $e+1$
\emph{інкапсуляція} KEM виконується до \emph{старого} публічного
Kyber-ключа, який входив до скомпрометованого стану. Свіжа Kyber-пара
$\mathit{pk\_new}_{e+1}$ лише \emph{генерується} на епосі $e+1$ і
\emph{передається} в заголовку, але першу безпечну інкапсуляцію до неї
контрагент може виконати тільки на епосі $e+2$. Отже,
постквантова гарантія повністю відновлюється на $e+2$ — звідси
\textbf{(б)}.

\emph{Мінімальність для цієї конструкції.} У \Aura{} без попередньо
опублікованого пулу ефемерних KEM-ключів ця затримка є мінімальною:
новий KEM-ключ спочатку генерується і публікується, а перша
інкапсуляція до нього можлива лише в наступному спрямованому кроці.
Отже, 2-крокова гібридна PCS є наслідком напрямковості KEM у цій
конструкції, а не дефектом реалізації. За компрометації лише одного
кінця перший крок може вже містити некомпрометований KEM-внесок peer.
\end{proof}

\begin{remark}
Порівняймо з PQXDH: там постквантовий внесок входить лише до
початкового X3DH. Після першого ж кроку ратчета весь похідний
матеріал залежить лише від класичного DH; отже, для трафіка,
збереженого опонентом, квантово-стійке PCS не відновлюється в межах
цієї схеми. \Aura{} відновлює його за 2~кроки. Це кількісно виражає
ефект KEM-ротації на кроках ратчета.
\end{remark}

\subsection{Теорема 5: конфіденційність та цілісність
повідомлень}\label{sec:thm5}

\begin{theorem}[Безпека повідомлень]\label{thm:conf-int}
За \MRAE{}-стійкості \AEAD{} та \PRF{}-стійкості ланцюгового KDF
кожний ключ повідомлення $\mk_{e,n}$ є псевдовипадковим, а ciphertext
забезпечує \INDCPA{}-конфіденційність і \INTCTXT{}-цілісність:
\begin{equation}\label{eq:msg-security-bound}
  \adv^{\mathsf{Conf}}_{\Aura}(\advA)
  \le
  \adv^{\eCK}_{\Aura}(\advB_0)
  + q_m \cdot \adv^{\PRF}_{\mathsf{chain}}(\advB_1)
  + q_m \cdot \adv^{\MRAE}_{\AEAD}(\advB_2)
  + \frac{q_m^2}{2^{96}}.
\end{equation}
\end{theorem}

\begin{proof}[Ескіз доведення]
Спершу кореневий ключ замінюється випадковим значенням за
\cref{thm:ake,thm:fs,thm:pcs}. Далі гібридним аргументом за індексами
ланцюга кожний $\mk_{e,n}$ замінюється випадковим ключем
(\PRF{}-стійкість chain-KDF). Після цього \MRAE{}-стійкість
\AEAD{} дає \INDCPA{} + \INTCTXT{} для корисного навантаження і
метаданих. Доданок $q_m^2/2^{96}$ консервативно обмежує колізії
96-бітного nonce-простору; всередині однієї сесії nonce унікальні за
побудовою через монотонні лічильники.
\end{proof}

\subsection{Теорема 6: стійкість до повторів}\label{sec:thm6}

\begin{theorem}[Антирепрейова стійкість]\label{thm:replay}
Для будь-якого опонента, що повторює не більш ніж $q_r$ повідомлень у
межах вікна $W = 2048$ nonce-значень:
\begin{equation}\label{eq:replay-bound}
  \adv^{\mathsf{Replay}}_{\Aura}(\advA)
  \le
  \frac{W \cdot q_r}{2^{96}}
  + \adv^{\INTCTXT}_{\AEAD}(\advB).
\end{equation}
\end{theorem}

\begin{proof}[Ескіз доведення]
Кожне повідомлення несе унікальну пару $\langle\epochv, N\rangle$.
Приймальний бік підтримує LRU-кеш до 2048~$\langle\epochv,
N\rangle$-пар; повтор у межах вікна виявляється точним збігом.
Повідомлення зі старішої епохи відхиляється, бо відповідний
ланцюговий ключ уже затерто (\texttt{zeroize}) і його не можна
відновити (одно-направленість HMAC-ратчета, \cref{thm:fs}).
Модифікований повтор потребує підробки AEAD-тегу і зводиться до
\INTCTXT{}-стійкості. Доданок $W q_r/2^{96}$ є консервативною межею
для випадкового збігу 96-бітного nonce у вікні дедуплікації.
\end{proof}

\subsection{Конкретні межі безпеки}\label{sec:bounds}

\begin{table}[H]
\centering
\caption{Конкретні рівні безпеки компонентів та композитні
межі.}\label{tab:bounds}
\small
\renewcommand{\arraystretch}{1.25}
\begin{tabular}{@{}llr@{}}
\toprule
Компонент / припущення & Стандарт & Біти \\
\midrule
X25519 / Gap-CDH & Curve25519 & 128 \\
ML-KEM-768 / Module-LWE & FIPS~203, category~3 & NIST~3 \\
AES-256-GCM-SIV / MRAE & RFC~8452 & 128 \\
HKDF-SHA256 / dual-PRF & RFC~5869 & PRF, 256-бітовий вихід \\
HMAC-SHA256 / PRF & FIPS~198-1 & PRF, 256-бітовий вихід \\
Ed25519 / SUF-CMA & EdDSA & 128 \\
\midrule
\multicolumn{3}{@{}l}{\textbf{Композитні межі протоколу}} \\
\midrule
Гібридний комбінатор (\cref{thm:combiner}), Gap-CDH $\lor$ ML-KEM & --- & 128 \\
AKE-стійкість (\cref{thm:ake}), усі припущення & --- & 128 \\
Пряма секретність (\cref{thm:fs}), Gap-CDH зі стертим DH-ефемером & --- & 128 \\
PCS класична (\cref{thm:pcs}а), Gap-CDH, 1~крок & --- & 128 \\
PCS гібридна (\cref{thm:pcs}б), Gap-CDH $+$ ML-KEM, 2~кроки & --- & 128 \\
Безпека повідомлень (\cref{thm:conf-int}), PRF $+$ MRAE & --- & 128 \\
\bottomrule
\end{tabular}
\end{table}

\paragraph{Втрата безпеки.}
Домінуючий фактор втрати — це $q_s$ (кількість сесій) у
\cref{thm:ake}, що є невід'ємною властивістю аргументу вгадування
моделі \eCK{}. Для типових розгортань із $q_s \le 2^{30}$ сесій
ефективна стійкість становить $128 - 30 = 98$~біт проти класичного
опонента. ML-KEM-768 трактується як NIST category~3; числові Core-SVP
оцінки треба подавати як орієнтовні та залежні від методики, а не як
сталий параметр протоколу.

\paragraph{Тісність.}
Редукції теорем~\ref{thm:fs} та~\ref{thm:pcs} є \emph{тісними} (без
множника вгадування поза гібридним комбінатором), бо опонент
націлюється на конкретну епоху, і редукція безпосередньо вкладає
челендж.

% ═════════════════════════════════════════════════════════════
% 3. ФОРМАЛЬНА ВЕРИФІКАЦІЯ
% ═════════════════════════════════════════════════════════════
\section{Машинно-перевірена формальна верифікація}\label{sec:formal}

Ігрові доведення (\cref{sec:theorems}) доповнено машинно-перевіреною
символьною верифікацією двома взаємодоповнювальними
інструментами: Tamarin~Prover~\cite{meier2013} для точного
trace-орієнтованого міркування та ProVerif~\cite{blanchet2016} для
аналізу на основі Horn-формул. У завершену машинно-доведену частину
входять 10~лем Tamarin і 3~з~6 ProVerif-зобов'язань; три ProVerif-запити
не враховано як машинно-доведені результати.

\subsection{Tamarin Prover}\label{sec:tamarin}

\paragraph{Модель рукостискання.}
Модель \texttt{aura\_handshake.spthy} використовує вбудовану в
Tamarin рівняльну теорію Діффі--Геллмана та власні функційні
символи для Kyber-768:
\begin{align}
  \texttt{kyber\_pk}(sk) \to kpk,\quad
  \texttt{kyber\_encaps}(ss, kpk) \to ct,\quad
  \texttt{kyber\_decaps}(ct, sk) = ss. \label{eq:tamarin-kyber-symbols}
\end{align}
Моделюються лише 2~з~4 X3DH DH-операцій (це коректна абстракція:
зменшення кількості DH лише \emph{послаблює} протокол, тобто дає
опоненту більше сили). Опонент Dolev--Yao має повний контроль над
мережею; два правила компрометації моделюють класичну (лише X25519)
та повну (X25519 $+$ Kyber) компрометацію.

\begin{table}[H]
\centering
\caption{Tamarin: модель рукостискання — 6/6 лем
підтверджено.}\label{tab:tamarin-hs}
\small
\begin{tabularx}{\textwidth}{@{}clX@{}}
\toprule
\# & Лема & Властивість \\
\midrule
1 & \texttt{session\_key\_secrecy} & Гібридний кореневий секрет безпечний, доки не скомпрометовано хоча б одну сторону \\
2 & \texttt{mutual\_authentication} & Двостороннє підтвердження ключа запобігає UKS \\
3 & \texttt{responder\_authentication} & Симетрично: коміт респондента $\Rightarrow$ ініціатор виконувався \\
4 & \texttt{forward\_secrecy\_hybrid} & Лише класична компрометація (без Kyber-SK) не руйнує секретність ключа \\
5 & \texttt{key\_confirmation} & Партнерські сесії деривують ідентичні кореневі ключі \\
6 & \texttt{session\_exists} & Досяжність: чесна сесія може завершитися \\
\bottomrule
\end{tabularx}
\end{table}

\paragraph{Модель ратчета.}
Модель \texttt{aura\_ratchet.spthy} абстрагує DH як класичний KEM
(щоб уникнути нетермінування Tamarin при DH у циклічних правилах
стану) та використовує термінальну (однокрокову) модель.

\begin{table}[H]
\centering
\caption{Tamarin: модель ратчета — 4/4 леми
підтверджено.}\label{tab:tamarin-rt}
\small
\begin{tabularx}{\textwidth}{@{}clX@{}}
\toprule
\# & Лема & Властивість \\
\midrule
1 & \texttt{pcs\_sender\_compromise} & Після компрометації стану відправника свіжий ратчет-ключ секретний, доки не скомпрометовано і приймача \\
2 & \texttt{ratchet\_key\_secrecy} & Ратчет-ключ секретний, доки не скомпрометовано обидві сторони \\
3 & \texttt{key\_agreement} & Обидві сторони деривують ідентичний кореневий ключ \\
4 & \texttt{ratchet\_exists} & Досяжність: чесний крок ратчета може завершитися \\
\bottomrule
\end{tabularx}
\end{table}

Лема \texttt{pcs\_sender\_compromise} перевіряє PCS-властивість у
термінальній однокроковій абстракції ратчета: навіть за повністю
скомпрометованого стану відправника свіжий гібридний крок ратчета
відновлює секретність, якщо одержувач не скомпрометований.

\subsection{ProVerif}\label{sec:proverif}

Активна модель \texttt{aura.pv} покриває handshake/message path:
моделює всі чотири X3DH DH-операції, перевірку Ed25519-підпису,
початковий chain KDF та доповнені (augmented) HKDF-$\mathit{info}$-рядки.
Використано один публічний канал Долева--Яо і 10~символьних типів, а
також необмежену реплікацію сесій. PCS-ратчет аналізується Tamarin-моделлю
ратчета та ігровою теоремою, а не активним ProVerif-процесом.
Запити Q1 і Q4 є перевірками секретності без компрометації у фазі~0.
Фазовий запит Q6 є stress query: фаза~0 — чесне виконання, фаза~1
розкриває всі довгострокові ключі.

\begin{table}[H]
\centering
\caption{ProVerif: 3/6 зобов'язань підтверджено.}\label{tab:proverif}
\small
\begin{tabularx}{\textwidth}{@{}clXc@{}}
\toprule
\# & Запит & Властивість & Статус \\
\midrule
Q1 & Секретність KEM-секрету & KEM-секрет фази~0, що живить $\rk_0$, прихований від опонента & \checkmark \\
Q2 & Автентифікація (неін'єкт.) & Завершення ініціатора $\Rightarrow$ респондент прийняв & \checkmark \\
Q3 & Автентифікація (ін'єктивна) & Ін'єктивна відповідність Q2 & н/з \\
Q4 & Секретність повідомлення & Зашифроване навантаження фази~0 секретне & \checkmark \\
Q5 & Цілісність повідомлення & Відповідність подій send/receive & $\times$ \\
Q6 & Пряма секретність (фазова) & Компрометація фази~1 не розкриває KEM-секрет фази~0 & $\times$ \\
\bottomrule
\end{tabularx}
\end{table}

\paragraph{Аналіз непідтверджених запитів.}
Q3, Q5 та Q6 не отримали машинного доведення в поточній ProVerif-моделі.
Q3 є зобов'язанням ін'єктивної автентифікації, яке не завершується в
стандартній чотири-DH моделі в межах артефактного бюджету. Q5 є занадто
широким непартнерським запитом цілісності: ProVerif знаходить активну
трасу, де опонент сам запускає сесію з респондентом і вводить власне
навантаження, тому така траса не відповідає теоремі для партнерських
сесій. Q6 питає секретність сирого KEM-секрету після розкриття
довгострокового KEM-SK; цей секрет справді відновлюється декapsulation-ом
і не є заявленою FS-властивістю. Автентифікацію, цілісність у партнерських
сесіях та PCS підтримано Tamarin-моделями, ігровими доведеннями і тестами
реалізації.

\paragraph{Взаємодоповнюваність інструментів.}
Tamarin дає точне DH-міркування, але потребує абстракцій проти
нетермінування; ProVerif опрацьовує необмежені сесії з усіма
чотирма DH-операціями і добре виявляє надто широкі формулювання запитів.
У підсумку доведено 10~лем Tamarin і 3~з~6 ProVerif-зобов'язань для
handshake/message-path моделі; Q3/Q5/Q6 не включено до
машинно-доведеної частини результату.

% ═════════════════════════════════════════════════════════════
% 4. РЕАЛІЗАЦІЯ
% ═════════════════════════════════════════════════════════════
\section{Реалізація}\label{sec:implementation}

Протокол \Aura{} реалізовано мовою Rust (edition~2021, MSRV~1.86).
Поточний зріз репозиторію містить близько 35{,}9~тис. рядків у
\texttt{src/}, 30{,}1~тис. рядків тестового коду та 1{,}1~тис. рядків
бенчмарків.

\subsection{Вибір криптографічних бібліотек}\label{sec:libs}

\begin{table}[H]
\centering
\caption{Криптографічні залежності реалізації.}\label{tab:libs}
\small
\begin{tabularx}{\textwidth}{@{}llX@{}}
\toprule
Примітив & Крейт & Призначення \\
\midrule
AES-256-GCM-SIV & \texttt{aes-gcm-siv} 0.11 & MRAE-AEAD обох рівнів \\
HKDF-SHA256 & \texttt{hkdf} 0.12, \texttt{sha2} & деривація ключів \\
HMAC-SHA256 & \texttt{hmac} 0.12 & ланцюговий ратчет, цілісність стану \\
X25519 & \texttt{x25519-dalek} & класичний DH \\
Ed25519 & \texttt{ed25519-dalek} & підписи бандлу \\
ML-KEM-768 & \texttt{ml-kem} 0.2 & постквантовий KEM (FIPS~203) \\
\texttt{zeroize} & \texttt{zeroize} & затирання чутливої пам'яті \\
\bottomrule
\end{tabularx}
\end{table}

\subsection{Захищене керування пам'яттю}\label{sec:secmem}

Довгострокові приватні DH/KEM/Ed25519 ключі зберігаються через
\texttt{SecureMemoryHandle}. Root/chain/metadata/message ключі є
байтовими буферами стану: вони явно затираються після використання, а
на диску захищаються sealed-state. \texttt{SecureMemoryHandle}:
\begin{itemize}
  \item викликає \texttt{mlock(2)} для запобігання вивантаженню
    сторінок із ключами у своп;
  \item затирає вміст через \texttt{zeroize} при звільненні
    (\texttt{Drop}), у тому числі на всіх шляхах помилок;
  \item проміжний матеріал ($\mathit{ikm}$, $\mathit{dh\_ss}$,
    $k_{\KEM}$, $\mathit{hybrid\_ikm}$) затирається явно одразу
    після використання, ще до завершення функції.
\end{itemize}
Прапорець \texttt{no-secure-memory} (лише для тестових середовищ)
вимикає \texttt{mlock}; його заборонено у продуктовій збірці.

\subsection{C FFI для крос-платформності}\label{sec:ffi}

Шар C~FFI (прапорець \texttt{ffi}, типи виходу
\texttt{staticlib}/\texttt{cdylib}) експортує протокол для Swift/iOS,
Android, C++ та .NET-клієнтів. Серверна (Rust) сторона використовує
крейт безпосередньо через модуль \texttt{api}, без FFI. FFI-шар
покрито 139~окремими тестовими сценаріями при збірці
\texttt{--features ffi}; усі помилки передаються кодами
(наприклад, \texttt{AuraErrorInvalidState}~$=16$), без розкручування
паніки через межу FFI.

\subsection{Тестове покриття}\label{sec:tests}

Реалізацію валідовано поточною матрицею з 750~тестових сценаріїв у
стандартній конфігурації та 889~сценаріїв при ввімкненому
\texttt{ffi}:

\begin{table}[H]
\centering
\caption{Структура тестового покриття.}\label{tab:tests}
\small
\begin{tabularx}{\textwidth}{@{}Xr@{}}
\toprule
Категорія & Сценаріїв \\
\midrule
Unit-тести основного крейту (\texttt{cargo test --lib}) & 36 \\
Регресійні тести публічного Rust API & 71 \\
Інтеграційні тести протоколу та стану сесій & 392 \\
Attack-PoC та зловмисні регресійні сценарії & 40 \\
Контракти заголовків, час, malformed-input та property-based сценарії & 13 \\
Додаткові real-time media регресії поза межами теорем & 194 \\
Додаткові поточні регресійні сценарії & 4 \\
\midrule
\textbf{Разом, стандартна конфігурація} & \textbf{750} \\
C~FFI при \texttt{cargo test --features ffi} & 139 \\
\textbf{Разом, feature-matrix із FFI} & \textbf{889} \\
\bottomrule
\end{tabularx}
\end{table}

Кожна з шести теорем безпеки відображена на емпіричні тести-валідатори
(test-to-proof mapping): наприклад, \cref{thm:pcs} валідовано тестами
\path|post_compromise_security_after_ratchet| та
\path|pcs_across_five_ratchet_steps|.

\subsection{Відповідність між моделлю, кодом і тестами}\label{sec:code-proof-map}

Формальні твердження цієї частини не замінюють інженерну валідацію:
вони задають, які властивості має мати будь-яка коректна реалізація.
Тому для \Aura{} важливо явно вказати, де саме абстрактні об'єкти
моделі з'являються в коді та якими тестами перевіряються їхні
реалізаційні межі. Така відповідність не є доказом сама по собі,
але вона зменшує ризик розриву між описаним у статті протоколом і його реалізацією.

\begin{table}[H]
\centering
\caption{Відповідність між доказовими об'єктами та реалізацією.}
\label{tab:code-proof-map}
\small
\begin{tabularx}{\textwidth}{@{}p{0.20\textwidth}YY@{}}
\toprule
Об'єкт моделі & Реалізаційний відповідник & Перевірка \\
\midrule
Гібридний X3DH &
\texttt{HandshakeInitiator}, \texttt{HandshakeResponder}; X25519 +
\texttt{KyberInterop} поверх \texttt{MlKem768} &
Tamarin \texttt{aura\_handshake.spthy}; API та integration-тести
рукостискання, key confirmation, tamper rejection \\
\addlinespace
Пере-ратчетний KEM-крок &
Ротація локальної DH/KEM-пари, новий ratchet header, augmented HKDF
із \texttt{Aura-Hybrid-Ratchet} &
Tamarin \texttt{aura\_ratchet.spthy}; тести PCS, direction-change,
out-of-order і chain exhaustion \\
\addlinespace
Дворівневе AEAD &
Окремі ключі \texttt{message key} і \texttt{metadata key}, обидва
через AES-256-GCM-SIV &
Теорема \ref{thm:conf-int}; тести metadata roundtrip, tampered
ciphertext/header, wrong key і nonce-boundary \\
\addlinespace
Replay-захист &
Bounded вікно для $\langle epoch, nonce\rangle$, кеш handshake-init
хешів, обмежені skipped-key структури &
Теорема \ref{thm:replay}; session replay, malformed-input і
state-recovery регресії в межах двостороннього протоколу \\
\addlinespace
Sealed state і rollback &
Двошарове AEAD-збереження, HMAC стану, зовнішній монотонний
лічильник і строге порівняння $<$ &
Тести sealed-state roundtrip, wrong-key, tamper, rollback і cold-start
equal-counter acceptance \\
\bottomrule
\end{tabularx}
\end{table}

\paragraph{Зріз, від якого рахуються цифри.}
Кількість тестових сценаріїв у \cref{tab:tests} отримано з
\texttt{cargo test -- --list} та
\texttt{cargo test --features ffi -- --list}; кількість рядків — із
підрахунку \texttt{*.rs} у \texttt{src/}, \texttt{tests/} та
\texttt{benches/}. Ці числа не входять до криптографічних припущень і
можуть змінюватися від коміту до коміту. Їхня роль у статті суто
відтворювальна: читач має бачити, який саме розмір реалізації та
тестової матриці відповідає наведеним твердженням.

\paragraph{Межі test-to-proof mapping.}
Тести не доводять \eCK{}-безпеку і не можуть замінити Tamarin чи
ігрові редукції. Їхня задача інша: перевірити, що реалізація не
порушує передумови моделі. Наприклад, теорема PCS припускає
атомарний ratchet-перехід; відповідні тести перевіряють rollback
після помилки, декілька послідовних ратчетів і неможливість
розшифрувати повідомлення старим станом. Теорема про replay-захист
припускає унікальність прийнятих nonce; тести перевіряють повтор
кадру, повтор групового повідомлення та відновлення replay-вікна
після sealed-state імпорту.

\paragraph{Що навмисно не переноситься у формальну модель.}
Формальні моделі не описують protobuf-кодування, allocation-пули
\texttt{SecureMemoryHandle}, FFI ownership-правила та конкретні
ліміти повідомлень. Ці частини перевіряються regression-тестами як
системні, а не криптографічні властивості. Такий поділ
зменшує розмір Tamarin/ProVerif моделей і залишає в них саме ті
аспекти, для яких потрібне символьне або ігрове міркування.

% ═════════════════════════════════════════════════════════════
% 5. ВІДТВОРЮВАНІСТЬ І МЕЖІ ТВЕРДЖЕНЬ
% ═════════════════════════════════════════════════════════════
\section{Відтворюваність і межі тверджень}\label{sec:reproducibility}

Криптографічна стаття про реалізований протокол має дві різні
відповідальності. Перша — сформулювати властивості протоколу в моделі,
де можна робити доведення. Друга — показати, що реалізація не відійшла
від цієї моделі у місцях, від яких залежать твердження. У цій роботі ці
дві відповідальності розділено явно: формальні моделі описують обмін
секретами, події автентифікації, ратчет і replay-властивості; Rust-тести
перевіряють кодування, межі буферів, FFI-контракти, sealed state,
обробку помилок і поведінку за межовими параметрами.

Таке розділення важливе практично. Якщо в коді змінити доменну мітку
\HKDF{}, розмір \KEM{}-ключа, правило прийняття sealed-state лічильника
або порядок оновлення send/recv-ланцюгів, це вже не косметична правка:
вона може змінити зв'язок між реалізацією і доказом. Натомість зміна
назви тесту, реорганізація модулів або додавання нової FFI-обгортки не
змінюють криптографічного твердження, якщо збережено ті самі інваріанти
протоколу. Саме тому числові характеристики реалізації в цій статті
подано як відтворювальний зріз, а не як частину припущень безпеки.

\subsection{Контур перевірки}

Мінімальний контур перевірки складається з чотирьох незалежних шарів:
переліку тестових сценаріїв, запуску тестів, формальної перевірки та
бенчмарків. Перелік тестів корисний для статті, бо дає стабільний
зріз покриття без залежності від часу виконання конкретної машини.
Повний запуск тестів потрібний для прийняття зміни в коді. Формальні
інструменти перевіряють властивості абстрактної моделі, а бенчмарки
показують вартість тих самих операцій у реалізації.

\begin{table}[tbp]
\centering
\caption{Команди відтворення та зміст перевірки.}
\label{tab:repro-commands}
\footnotesize
\begin{tabularx}{\textwidth}{@{}p{0.25\textwidth}YY@{}}
\toprule
Шар & Команда або артефакт & Що саме перевіряється \\
\midrule
Єдиний artifact-entrypoint &
\texttt{artifact quick mode} &
Фіксує версії інструментів, запускає paper vectors, deterministic
handshake/envelope/cross-ratchet/multi-epoch KATs і attack-PoC регресії; повні режими
\texttt{test}, \texttt{formal}, \texttt{paper}, \texttt{bench},
\texttt{full} описано в artifact-документі. \\
\addlinespace
Фіксовані vectors &
\texttt{test-vectors paper\_vectors} &
Стабільні вектори для HKDF, AES-256-GCM-SIV, seeded ML-KEM public key
і master-key identity derivation; deterministic
handshake/envelope/cross-ratchet/multi-epoch KATs перевіряють байти
\texttt{HandshakeInit}, \texttt{HandshakeAck}, session id, перший
post-handshake envelope, перший DH+KEM ratchet envelope і delayed delivery
через кілька ratchet-епох під test-only entropy feature. \\
\addlinespace
Зріз тестів за замовчуванням &
\texttt{cargo test -- --list} &
Кількість сценаріїв без FFI-прапорця; у поточному зрізі — 750 тестових
входів, включно з unit, API, integration, attack-PoC,
header/time/malformed/property тестами та додатковими out-of-scope
регресіями. \\
\addlinespace
Зріз тестів з FFI &
\texttt{cargo test --features ffi -- --list} &
Той самий набір плюс FFI-контракти, malformed-input сценарії та
перевірки C API; у поточному зрізі — 889 тестових входів. \\
\addlinespace
Розмір реалізації &
\texttt{find src tests benches -name '*.rs' -print} &
Підрахунок рядків Rust-коду: 35{,}9~тис. у \texttt{src/},
30{,}1~тис. у \texttt{tests/} і 1{,}1~тис. у
\texttt{benches/protocol\_bench.rs}. \\
\addlinespace
Tamarin &
\texttt{make handshake}; \texttt{make ratchet} у \texttt{formal/} &
Шість лем для гібридного X3DH та чотири леми для ратчета. Повна
комбінована модель зберігається як reference-модель, але не є
основним шляхом перевірки через складність DH-еквацій. \\
\addlinespace
ProVerif &
\texttt{make proverif} у \texttt{formal/} &
Чотири доведені запити для handshake/message-path моделі; два
недоведені запити зберігаються як ProVerif-зобов'язання, а не як
позитивні твердження статті. \\
\addlinespace
Бенчмарки &
\texttt{cargo bench --bench protocol\_bench} &
Латентність рукостискання, ratchet-кроку, шифрування/розшифрування,
sealed-state export/import і фіксованих traffic-profile сценаріїв. \\
\bottomrule
\end{tabularx}
\end{table}

Команди в \cref{tab:repro-commands} мають різний статус. Підрахунок
рядків і список тестів фіксують масштаб артефакта; вони самі нічого не
доводять. Tamarin і ProVerif перевіряють властивості моделі, але не
бачать деталей пам'яті, protobuf і FFI. Бенчмарки не є доказом
безпеки; вони обмежують твердження про практичність, бо PCS-профіль
оцінюється разом із вартістю ratchet-кроку.

\subsection{Декомпозиція формальної моделі}

Повна модель \texttt{formal/tamarin/aura.spthy} корисна як специфікація
подій і зв'язків між фазами, але для автоматичного доведення вона
занадто важка: DH-екваційна теорія швидко роздуває простір джерел.
Тому перевірку розкладено на \texttt{aura\_handshake.spthy} і
\texttt{aura\_ratchet.spthy}. Це не змінює змісту тверджень, якщо
межа між моделями збігається з реальною межею в коді: результатом
рукостискання є початковий root key, підтверджені ідентичності та
початкові send/recv-ланцюги; входом ратчета є саме такий стан і подія
отримання нового ratchet header.

Така декомпозиція також обмежує формулювання результатів. У ній не можна
стверджувати, що ProVerif доводить усі властивості повідомлень або
PCS-ратчета, якщо три запити залишаються недоведеними. Не можна також переносити леми
VoIP-моделі на основний двосторонній протокол без окремого формального
зв'язку. Тому в основний результат винесено рівно ті формальні факти,
які мають завершену перевірку: 6 лем Tamarin для handshake, 4 леми
Tamarin для ratchet і 3 доведені ProVerif-зобов'язання handshake/message-path
моделі. Решта артефактів
позначають напрям розширення, але не збільшують заявлений доказовий
обсяг.

\subsection{Зовнішні припущення реалізації}

Доведення не усувають потребу в зовнішніх припущеннях. Ідентичності
мають бути прив'язані до користувачів через PKI, TOFU, safety numbers
або інший канал автентифікації; протокол не може сам відрізнити новий
чесний пристрій від підміни каталогу ключів. Генератор випадковості
має надавати достатню ентропію для DH, ML-KEM, nonce-prefix і
sealed-state ключів. Монотонний лічильник sealed state має зберігатися
поза самим зашифрованим blob; інакше rollback-захист перетворюється на
перевірку цілісності blob, але не на захист від відновлення старої
копії сховища.

Окремо, \Aura{} не заявляє постквантову стійкість підписів: ідентичності
та підписані pre-key-и в поточній реалізації спираються на Ed25519.
Квантовий супротивник із CRQC може атакувати автентичність початкового
бандлу, навіть якщо KEM-компонент лишається стійким. Так само модель
двосторонньої сесії не є повною multi-device моделлю: fan-out сесій,
Sesame-подібна маршрутизація та політики device linking мають
аналізуватися як окремий шар.

Пам'ять із секретами захищається best-effort засобами ОС:
\texttt{mlock}, zeroize та, на Linux, \texttt{MADV\_DONTDUMP}. Це
зменшує ризик витоку через swap або core dump, але не є моделлю
захисту від повністю привілейованого локального супротивника.
Аналогічно, FFI-шар перевіряє довжини, null-вказівники та ownership
результатів, однак коректність життєвого циклу об'єктів у зовнішній
мові залишається відповідальністю інтегратора. Ці припущення не
послаблюють криптографічний результат; вони задають межу, за якою
починається системна безпека застосунку.

\subsection{Критерій прийняття майбутніх змін}

Для подальшого розвитку протоколу корисний простий критерій: будь-яка
зміна, що торкається секретного входу \HKDF{}, AAD, ratchet header,
nonce, epoch, sealed-state counter або replay-вікна, має супроводжуватися
оновленням принаймні одного з трьох артефактів — доказу, формальної
моделі або тесту, який явно фіксує новий інваріант. Якщо змінюється
лише API-обгортка або назва історичного типу на кшталт
\texttt{KyberInterop}, достатньо regression-тестів сумісності. Такий
критерій не замінює рецензування, але робить видимим головний ризик:
поступове розходження між протоколом, кодом і текстом.

% ═════════════════════════════════════════════════════════════
% 6. НЕГАТИВНІ СЦЕНАРІЇ ТА ОПЕРАЦІЙНІ ІНВАРІАНТИ
% ═════════════════════════════════════════════════════════════
\section{Негативні сценарії та операційні інваріанти}
\label{sec:negative-scenarios}

Окрема частина коректності протоколу проявляється не в успішному
обміні повідомленнями, а в тому, як реалізація поводиться на межі:
при пошкодженому input, повторній доставці, відновленні старого стану,
частковому crash між ratchet-кроком і записом на диск або при
підміні публічного бандлу сервером. Саме ці сценарії часто не видно з
основного опису протоколу, бо формальна модель абстрагує транспорт,
серіалізацію й локальне сховище. Для \Aura{} вони винесені в явні
операційні інваріанти: повідомлення або приймається з повним
оновленням стану, або відкидається без часткового зміщення epoch,
chain-key, replay-вікна чи sealed-state counter.

\paragraph{Пошкоджені та неповні кадри.}
Перший клас негативних сценаріїв — syntactically invalid input:
обрізані protobuf-поля, неправильні довжини X25519/Kyber-буферів,
невідомі версії envelope, відсутній ratchet header, завеликий payload
або некоректна AEAD-мітка. У криптографічній моделі такі повідомлення
зазвичай не породжують подій, але в коді вони є критичними, бо саме
там виникають помилки пам'яті, неоднозначні стани й різні канали
спостереження. Реалізаційне правило просте: усі перевірки формату
виконуються до оновлення секретного стану; AEAD failure не розкриває,
яка саме внутрішня перевірка не пройшла; помилки FFI повертаються як
коди стану без передачі частково ініціалізованих буферів назовні.

Цей інваріант потрібен і для доказової частини. Якщо malformed frame
міг би змінити лічильник повідомлень або просунути replay-вікно, тоді
опонент отримав би поведінку, якої немає в \cref{sec:model}. Якщо ж
помилка обробляється як чистий reject без side effect, її можна
моделювати як відсутність прийнятої події. Тому тести на malformed
input не є другорядними інженерними перевірками: вони захищають відповідність між
системною реалізацією та абстрактним поняттям прийняття повідомлення.

\paragraph{Доставка поза порядком.}
Асинхронний месенджер не має права припускати FIFO-транспорт.
Повідомлення можуть приходити із затримкою, дублюватися, втрачатися
або доставлятися після ratchet-переходу. Для класичного Double Ratchet
це розв'язується кешем пропущених message keys; у \Aura{} додається
ще один вимір — epoch, у якій був зафіксований metadata-key та
гібридний KEM-внесок. Відповідно, ключі пропущених повідомлень мають
бути обмежені не лише кількістю, а й належністю до конкретного
ratchet-контексту. Прийняття старого повідомлення не повинно
відкочувати поточний root key, а прийняття нового ratchet header не
повинно робити старі кешовані ключі валідними для іншої епохи.

Практична політика тут свідомо консервативна: implementation може
підтримувати bounded out-of-order window, але не повинна намагатися
``лікувати'' довільно розсинхронізовані сесії шляхом пошуку старих
станів. Якщо сторони втратили спільний контекст за межами replay-вікна,
коректна відповідь — нове рукостискання або протокольна ресинхронізація
вищого рівня. Це зберігає властивість, потрібну для PCS: минулий
компрометований стан не стає знову активним лише тому, що transport
доставив старий пакет.

\paragraph{Crash consistency.}
Найнебезпечніший практичний збій — аварійне завершення між
розшифруванням повідомлення і записом оновленого стану. Якщо plaintext
уже віддано застосунку, але replay-вікно або chain key не збережені,
повтор того самого ciphertext після перезапуску може бути прийнятий
удруге. Тому прийняття повідомлення має трактуватися як транзакція:
обчислення нового стану, перевірка AEAD, оновлення replay-вікна,
sealed-state export і тільки після цього повернення успішного результату
виклику. У середовищах, де застосунок не може гарантувати атомарний
запис, потрібен зовнішній journal або write-ahead marker.

Це не змінює криптографічних рівнянь, але змінює систему безпеки.
Формальна теорема про replay-захист припускає, що подія прийняття
однозначно записана в стані сторони. Якщо реалізація допускає crash
після прийняття без durable update, то порушується не AEAD і не HKDF,
а саме передумова про стан. Тому sealed state у цій статті розглянуто
разом із rollback-лічильником і FFI-контрактами: без них доказ
правильний для протоколу, але недостатній для реального клієнта.

\paragraph{Каталог ключів і downgrade.}
Сервер, що роздає пре-ключові бандли, не є довіреним джерелом
автентичних ключів. Він може затримувати оновлення, повертати
старий signed pre-key, приховувати наявність KEM-ключа або змушувати
сторони домовлятися про слабший профіль сумісності. \Aura{} не
вирішує проблему глобальної прозорості каталогу; вона задає локальне
правило: якщо peer заявлений як такий, що підтримує гібридний профіль,
відсутність KEM-компонента або зміна параметрів має бути видимою
подією policy failure, а не мовчазним fallback.

У продуктовій системі це правило зазвичай доповнюється журналом
версій ключів, safety-number перевіркою, key transparency або іншим
механізмом виявлення підміни. Для статті важливо інше: downgrade не
можна ховати під тим самим security claim, що й повний гібридний
режим. Якщо KEM-компонент відсутній, твердження про
harvest-now-decrypt-later-захист
для цього рукостискання не застосовується. Така явна межа робить
протокол чесним: він може мати сумісні режими, але кожен режим має
окремий набір властивостей.

\paragraph{Версіонування.}
Реальна бібліотека живе довше за одну статтю, тому формат envelope,
домени \HKDF{}, AAD і параметри KEM мають версіонуватися так, щоб
старий ciphertext не набував нового змісту після оновлення коду.
Версія протоколу входить до автентифікованого контексту; зміна
семантики поля вимагає нової версії, а не неявної інтерпретації того
самого byte layout. Це особливо важливо для metadata encryption:
заголовок зашифрований, але його контекст усе одно має бути
однозначним для сторін, які виконують розшифрування.

З практичного погляду, версіонування зменшує ризик ``тихих''
несумісностей між мобільними клієнтами, desktop-клієнтами й серверними
компонентами. З наукового погляду, воно фіксує об'єкт доведення:
твердження цієї статті стосуються конкретного набору доменів,
параметрів і правил обробки стану. Майбутні зміни можуть бути
сумісними з цим набором або замінювати його, але не повинні непомітно
змішувати два різні протоколи під однією назвою.

% ═════════════════════════════════════════════════════════════
% 7. ЕКСПЕРИМЕНТАЛЬНА ОЦІНКА
% ═════════════════════════════════════════════════════════════
\section{Експериментальна оцінка}\label{sec:evaluation}

Бенчмарки є single-host вимірюваннями Criterion.rs з
\texttt{cargo bench --bench protocol\_bench}: 200~семплів і 5~с
вимірювання. Значення залежать від hardware і є зрізом вартості
реалізації, а не припущенням безпеки. Artifact фіксує OS/CPU, версії
Rust, build-profile context, SHA-256 \texttt{Cargo.lock}, stdout-логи та
Criterion result files.
Публічний benchmark-harness також містить фіксовану групу
\texttt{pq\_traffic\_profiles}: 128 повідомлень по
256~байт для порівняння handshake-only PQ, sparse chain-boundary PQ, periodic
chain-boundary PQ та per-reply PQ rekeying в одному вимірювальному контурі.

\begin{table}[H]
\centering
\caption{Латентності основних операцій.}\label{tab:bench-core}
\small
\begin{tabular}{@{}lr@{}}
\toprule
Операція & Латентність \\
\midrule
Створення ідентичності (5~OPK) & $\sim$1{,}5~мс \\
Повне рукостискання (keygen $+$ X3DH $+$ Kyber $+$ confirm) & $\sim$1{,}1~мс \\
Рукостискання (з готовою ідентичністю) & $\sim$0{,}7~мс \\
Гібридний крок ратчета (X25519 $+$ Kyber-768) & $\sim$259~мкс \\
Шифрування 256~Б (той самий ланцюг, без ратчета) & $\sim$17~мкс \\
Розшифрування 256~Б (той самий ланцюг) & $\sim$21~мкс \\
\midrule
Kyber-768 keygen / encaps / decaps & $\sim$80 / 100 / 90~мкс \\
X25519 скалярне множення & $\sim$50~мкс \\
HKDF-SHA256 (32~Б виходу) & $\sim$2~мкс \\
Експорт / імпорт запечатаного стану & $\sim$30 / 35~мкс \\
\bottomrule
\end{tabular}
\end{table}

\begin{table}[H]
\centering
\caption{Decrypt latency для попередньо згенерованих envelope за розміром payload.}\label{tab:bench-payload}
\small
\begin{tabular}{@{}rr@{}}
\toprule
Payload & Decrypt latency \\
\midrule
64~Б & $\sim$35~мкс \\
256~Б & $\sim$38~мкс \\
1~KiB & $\sim$45~мкс \\
4~KiB & $\sim$60~мкс \\
16~KiB & $\sim$120~мкс \\
64~KiB & $\sim$350~мкс \\
256~KiB & $\sim$1{,}2~мс \\
1~MiB & $\sim$4{,}5~мс \\
\bottomrule
\end{tabular}
\end{table}

\begin{table}[H]
\centering
\caption{Накладні витрати ratchet-кроку.}\label{tab:bench-overhead}
\small
\begin{tabular}{@{}lr@{}}
\toprule
Компонент & Латентність \\
\midrule
X25519 DH scalar multiplication (класична база) & $\sim$50~мкс \\
Kyber-768 encap $+$ decap & $\sim$190~мкс \\
Повний гібридний ratchet step & $\sim$259~мкс \\
\midrule
Alternating throughput (ratchet на кожному 256~Б повідомленні) & $\sim$280~мкс/msg \\
Burst decrypt throughput (pre-generated envelopes, 256~Б) & $\sim$15~мкс/msg \\
\bottomrule
\end{tabular}
\end{table}

\paragraph{Аналіз накладних витрат.}
У цих вимірюваннях повний гібридний ratchet step коштує близько
$259$~мкс проти близько $50$~мкс для X25519. Pre-generated burst decrypt
для повідомлень в одному напрямку вимірюється на рівні
$\sim$15~мкс/msg; alternating-профіль із 256~Б payload, де encryption і
decryption входять у вимірюваний цикл, вимірюється на рівні
$\sim$280~мкс/msg. Повне рукостискання на benchmark-host займає близько
$1{,}1$~мс.

\paragraph{Накладні витрати на трафік.}
Кожен крок ратчета додає $1184 + 1088 = 2272$~байти (публічний
Kyber-ключ $+$ Kyber-шифротекст) плюс 32~байти нового X25519-ключа.
Для повідомлень в одному напрямку (без ратчета) накладні витрати на
заголовок зводяться до зашифрованих метаданих та номера епохи.

\section{Порівняння з актуальними постквантовими протоколами}\label{sec:comparison}

Коректна точка відліку для 2026~року ширша за \PQXDH{}.
Handshake-only постквантовий захист уже не є верхньою межею для
месенджерних протоколів: Apple \PQThree{} та \Signal{} \SPQR{}/Triple
Ratchet вводять постквантовий матеріал у подальшу еволюцію стану.
Тому порівняння нижче розділяє три питання: де саме з'являється
постквантова ентропія, як швидко система відновлюється після
компрометації та який рівень публічного підтвердження має конструкція.

\begin{table}[H]
\centering
\caption{Позиціювання \Aura{} відносно сучасних месенджерних і
групових протоколів.}\label{tab:sota-comparison}
\scriptsize
\begin{tabularx}{\textwidth}{@{}p{0.16\textwidth}p{0.19\textwidth}p{0.18\textwidth}p{0.17\textwidth}X@{}}
\toprule
Підхід & Де використовується PQ-матеріал & PCS після компрометації & Публічне підтвердження & Висновок для цієї роботи \\
\midrule
\Signal{} X3DH/Double Ratchet~\cite{signal-x3dh,signal-double-ratchet} &
Не використовується &
Класичне PCS через DH-ратчет &
Відкрита специфікація; академічні аналізи &
Базова модель FS/PCS, але без квантової стійкості. \\
\addlinespace
\Signal{} \PQXDH{}~\cite{signal-pqxdh} &
Початкове рукостискання &
Після handshake PQ-внесок не оновлюється &
Відкрита специфікація &
Захищає від harvest-now-decrypt-later для старту сесії, але не задає PQ-PCS для довгого каналу. \\
\addlinespace
\Signal{} \SPQR{}/Triple Ratchet~\cite{signal-spqr,signal-double-ratchet,mlkem-braid} &
Sparse PQ-ратчет поруч із класичним Double Ratchet &
PQ-PCS досягається через SCKA-епохи; прогрес залежить від доставки PQ-частин &
Відкриті специфікації; формальна верифікація як частина процесу розробки Signal &
Ключовий відкритий референс у сім'ї Signal; \Aura{} треба порівнювати саме з ним, а не лише з \PQXDH{}. \\
\addlinespace
Apple \PQThree{}~\cite{apple-pq3,stebila-pq3,linker-pq3} &
Початкове встановлення та періодичне in-band пере-ключення &
Заявлене і проаналізоване FS/PCS проти класичних і квантових опонентів &
Публічний опис, ігровий аналіз, Tamarin-аналіз; реалізація не є відкритим протокольним артефактом &
Промисловий baseline; \Aura{} не повинна заявляти першість у самій ідеї PQ-ратчета. \\
\addlinespace
\MLS{} RFC~9420~\cite{rfc9420} &
У базовому стандарті не постквантовий &
Епохальне PCS для груп через tree updates &
IETF-стандарт &
Референс для групової еволюції стану, але не прямий конкурент двосторонньому каналу. \\
\addlinespace
\textbf{\Aura{}} &
Гібридний handshake та KEM-ротація на межах ратчета &
1 класичний крок; 2 гібридні кроки для PQ-відновлення &
Відкрита Rust-реалізація, 6 теорем, Tamarin/ProVerif, proof-to-code таблиця &
Внесок статті: компактна відкрита конструкція з явним зв'язуванням KEM-ключа, окремим захистом метаданих і перевірюваними інваріантами реалізації. \\
\bottomrule
\end{tabularx}
\end{table}

\begin{table}[H]
\centering
\caption{Стисла матриця відмінностей \Aura{} від найближчих
орієнтирів.}\label{tab:comparison}
\small
\begin{tabularx}{\textwidth}{@{}p{0.22\textwidth}YYYY@{}}
\toprule
Характеристика & \PQXDH{} & \SPQR{} & \PQThree{} & \textbf{\Aura{}} \\
\midrule
PQ-внесок після handshake & Ні & Так & Так & \textbf{Так} \\
Модель оновлення PQ-матеріалу & Одноразова KEM-інкапсуляція & Sparse PQ-ратчет & Періодичне in-band пере-ключення & \textbf{KEM-ротація на ратчет-межі} \\
Зв'язування нового KEM-ключа з transcript & KEM лише в handshake & Публічний SCKA/Braid state задає оновлення KEM-ключів & Публічні матеріали описують in-band PQ rekeying & \textbf{KEM public key включено в HKDF info} \\
Окремий ротаційний захист заголовка & Ні & У публічній специфікації не заявлено окремий envelope-metadata рівень & У публічному дизайні не заявлено незалежний per-epoch envelope key & \textbf{Окремий metadata-key епохи} \\
Відкрита реалізація як артефакт статті & Специфікація & Специфікація & Публічний дизайн, закрита реалізація & \textbf{Rust-код, тести, proof-to-code mapping} \\
Формалізована затримка PQ-PCS & Ні & Аналіз SCKA-епох & Опублікований PQ3-аналіз & \textbf{Теорема про асиметрію 1/2 кроки} \\
\bottomrule
\end{tabularx}
\end{table}

Порівняльний висновок такий. \Aura{} не конкурує з Apple \PQThree{}
за масштабом розгортання і не перевершує \Signal{} \SPQR{} самим фактом
наявності постквантового ратчета. Основний внесок цієї роботи полягає
в описі відтворюваної академічної конструкції, у якій новий
KEM-матеріал формально входить у стан каналу, KEM-публічний ключ
зв'язується з transcript без додаткового підпису, метадані мають власну
ротацію ключа, а кожне твердження безпеки прив'язане до конкретних
модулів реалізації.

Для подання статті залишаються три зовнішні обмеження, які не слід
маскувати як уже завершені результати: незалежна криптографічна
рецензія моделі та коду, release-tagged artifact із зафіксованими
версіями інструментів і незалежно відтворені baseline-вимірювання для
\PQXDH{}/\SPQR{}/\PQThree{}-подібних профілів. Вони посилюють
емпіричну частину роботи, але не змінюють сформульовані тут теореми.

% ═════════════════════════════════════════════════════════════
% 6. ДИСКУСІЯ
% ═════════════════════════════════════════════════════════════
\section{Дискусія та обмеження}\label{sec:discussion}

\paragraph{Накладні витрати на трафік.}
KEM-ротація на ратчет-межі має ціну $\sim$2{,}3~КБ на відповідний
ratchet-крок. Для текстового месенджера придатність залежить від
частоти зміни напрямку та latency/bandwidth budget продукту; для
сценаріїв «пінг-понг» (постійне чергування) витрати можуть домінувати.
Режим з ратчетом раз на $k$ повідомлень слід аналізувати як окремий
профіль із власним PCS/cost твердженням.

\paragraph{2-крокова гібридна PCS.}
Як доведено у \cref{thm:pcs}, двокрокова затримка квантового
PCS-відновлення є структурним наслідком KEM-public-key ratchet
конструкції \Aura{} без наперед опублікованих ефемерних KEM-ключів.
Ми припускаємо, що протокол із 1-кроковою гібридною PCS
для обох сторін одночасно потребував би або передавання свіжих KEM-ключів
обома сторонами в одному протокольному повідомленні, або постквантового
інтерактивного обміну ключами, аналогічного DH за роллю.

\paragraph{Обмеження ProVerif.}
Q3/Q5/Q6 не доводяться у ProVerif; Q3 не завершується в дефолтній
чотири-DH моделі, Q5 є хибним для непартнерської активної траси, а Q6
є хибним для сирого KEM-секрету після розкриття KEM-SK. Ці запити не
враховано як машинно-доведені результати; заявлені властивості спираються
на Tamarin-моделі, ігрові доведення та тести реалізації.

\paragraph{Термінологія Kyber/ML-KEM.}
У коді історична назва \texttt{KyberInterop} збережена як API-шар, але
криптографічна реалізація вже використовує крейт \texttt{ml-kem} і тип
\texttt{MlKem768}. Перехід на сертифікований FIPS-модуль ML-KEM у
продуктовому середовищі не змінює структури протоколу та доведень.

\paragraph{Групове шифрування.}
Поточна стаття розглядає двосторонній сценарій. Крейт містить
експериментальні модулі групового шифрування (TreeKEM-подібна
структура); їх формальний аналіз винесено за межі цієї роботи.

\paragraph{Профілі практичного використання.}
Один і той самий криптографічний механізм має різну ціну в різних
режимах застосування. Для асинхронного текстового месенджера природним
є режим повного гібридного ратчета при кожній зміні напрямку: у
виміряному асинхронному текстовому профілі латентність ratchet-кроку
мала відносно людської затримки відповіді, а bandwidth-overhead
залежить від частоти зміни напрямку.
У такому профілі перевага безпеки проста: після компрометації стану
перший чесний крок повертає класичну PCS, а другий гібридний крок
повертає внесок KEM.

Для низьколатентних каналів, зокрема голосових або відео-сесій, повна
ротація KEM на кожному короткому кадрі була б неправильним інженерним
вибором. У цьому режимі базовий двосторонній handshake доцільно
використовувати для автентифікованого встановлення початкового секрету,
а частоту медіа-ратчета задавати окремим параметром часу або кількості
кадрів. Такий режим не змінює доведень для основного протоколу, але й
не повинен автоматично успадковувати всі їхні твердження: медіа-шар має
власні припущення про втрати пакетів, jitter, rekey-ack і життєвий цикл
дзвінка.

Для довгоживучих desktop- або server-side сесій ключовим стає sealed
state. Тут найбільший ризик не в одному невдалому ciphertext, а в
повторному імпорті старої копії стану після backup restore, snapshot або
міграції сховища. Тому монотонний лічильник sealed state треба
розглядати як частину протокольної інтеграції, а не як локальну
оптимізацію серіалізації. Без зовнішнього водяного знака HMAC і AEAD
захищають цілісність blob, але не можуть відрізнити свіжий blob від
старого коректного blob.

Ці профілі показують межу між протоколом і продуктом. Стаття доводить
властивості двостороннього обміну повідомленнями з конкретними
інваріантами ратчета, nonce, replay-вікон і sealed-state counter.
Вибір частоти ратчета, політики backup, синхронізації кількох пристроїв
і параметрів VoIP-кадрів є рівнем системного дизайну. Він має спиратися
на ці інваріанти, але не повинен змішуватися з ними в одному
неперевіреному твердженні.

\section{Висновки}\label{sec:conclusion}

Розглянуто \Aura{} — гібридний постквантовий протокол захищеного
обміну повідомленнями. Ключовий технічний результат полягає в тому, що
постквантовий компонент перенесено з одноразового рукостискання в
механізм станової еволюції сесії. \Aura{} створює гібридний початковий
секрет і регулярно повертає KEM-внесок у кореневий ключ після
компрометації стану. Це відрізняє протокол від підходів, де
постквантовість завершується на етапі X3DH.

Конкретний внесок роботи складається з п'яти перевірюваних елементів:
\emph{(Н1)}~напрямковий гібридний Double Ratchet із KEM-внеском на
ratchet-межах; \emph{(Н2)}~augmented HKDF, де новий KEM-публічний ключ
входить до параметра $\mathit{info}$ без додаткового підпису;
\emph{(Н3)}~дворівневе AEAD з ротаційним ключем метаданих;
\emph{(Н4)}~формалізація межі PCS-відновлення
«1~класичний крок / 2~гібридні кроки» для цієї конструкції;
\emph{(Н5)}~ланцюг перевірки, що поєднує 6~теорем ігрових редукцій,
10~лем Tamarin, 3~доведені ProVerif-зобов'язання handshake/message-path моделі і відображення доказових
інваріантів на Rust-тести.

Гібридний комбінатор досягає гарантії «\emph{або-або}» — безпеки за
умови стійкості \emph{принаймні одного} з X25519 чи Kyber-768. Усі
шість теорем доведено редукціями до стандартних припущень. Номінальна
класична міцність компонентів становить 128~біт; з урахуванням
множника $q_s$ в AKE-редукції для $q_s\le 2^{30}$ ефективна межа
становить 98~біт. Для ML-KEM-768 використовується категорія
NIST category~3; конкретні Core-SVP бітові оцінки залежать від методики.
Реалізація мовою Rust
(35{,}9~тис. рядків у \texttt{src/}; повний репозиторний зріз має
750/889~тестових сценаріїв залежно
від прапорця \texttt{ffi}) фіксує відтворюваний зріз продуктивності для
оцінювання протоколу в месенджерних сценаріях.

\paragraph{Напрями подальших досліджень.}
\emph{(а)}~налаштовуваний компроміс частоти ратчета проти
накладних витрат на трафік; \emph{(б)}~конструкції з однокроковою
гібридною PCS для обох сторін, наприклад через одночасну передачу
свіжих KEM-ключів або постквантовий інтерактивний обмін ключами;
\emph{(в)}~розширення формального аналізу на групове шифрування;
\emph{(г)}~валідація сертифікованого FIPS-модуля ML-KEM у продуктовому
ланцюгу постачання.

\section*{Подяки}

Автор дякує OpenAI Codex за допомогу з перевіркою узгодженості коду й
документації, LaTeX-діагностикою та редакційним вичитуванням.
Відповідальність за протокол, твердження безпеки і остаточний текст
належить автору.

% ═════════════════════════════════════════════════════════════


% ═════════════════════════════════════════════════════════════
\begin{thebibliography}{99}

\bibitem{signal-x3dh}
M.~Marlinspike, T.~Perrin.
\newblock The X3DH Key Agreement Protocol, Revision~1.
\newblock Signal Foundation, Specification, 2016.
\newblock \url{https://signal.org/docs/specifications/x3dh/}

\bibitem{cohn-gordon2020}
K.~Cohn-Gordon, C.~Cremers, B.~Dowling, L.~Garratt, D.~Stebila.
\newblock A formal security analysis of the Signal messaging protocol.
\newblock \emph{Journal of Cryptology}, 33(4):1914--1983, 2020.
\newblock \url{https://doi.org/10.1007/s00145-020-09360-1}

\bibitem{signal-pqxdh}
E.~Kret, R.~Schmidt.
\newblock The PQXDH Key Agreement Protocol, Revision~3.
\newblock Signal Foundation, Specification, 2024.
\newblock \url{https://signal.org/docs/specifications/pqxdh/}

\bibitem{signal-double-ratchet}
T.~Perrin, M.~Marlinspike, R.~Schmidt.
\newblock The Double Ratchet Algorithm, Revision~4.
\newblock Signal Foundation, Specification, 2025.
\newblock \url{https://signal.org/docs/specifications/doubleratchet/}

\bibitem{signal-spqr}
G.~Connell, R.~Schmidt.
\newblock Signal Protocol and Post-Quantum Ratchets.
\newblock Signal Blog, 2025.
\newblock \url{https://signal.org/blog/spqr/}

\bibitem{mlkem-braid}
R.~Schmidt.
\newblock The ML-KEM Braid Protocol.
\newblock Signal Foundation, Specification, Revision~1, 2025.
\newblock \url{https://signal.org/docs/specifications/mlkembraid/}

\bibitem{apple-pq3}
Apple Security Engineering and Architecture.
\newblock iMessage with PQ3: The new state of the art in quantum-secure
  messaging at scale.
\newblock Apple Security Research, 2024.
\newblock \url{https://security.apple.com/blog/imessage-pq3/}

\bibitem{stebila-pq3}
D.~Stebila.
\newblock Security analysis of the iMessage PQ3 protocol.
\newblock Technical report, 2024.
\newblock \url{https://www.douglas.stebila.ca/research/papers/Apple-Stebila24/}

\bibitem{linker-pq3}
F.~Linker, R.~Sasse, D.~Basin.
\newblock A Formal Analysis of Apple's iMessage PQ3 Protocol.
\newblock In \emph{USENIX Security 2025}, pp.~5015--5034.
\newblock \url{https://www.usenix.org/conference/usenixsecurity25/presentation/linker}

\bibitem{brendel2020}
J.~Brendel, M.~Fischlin, F.~G{\"u}nther, C.~Janson, D.~Stebila.
\newblock Towards Post-Quantum Security for Signal's X3DH Handshake.
\newblock In \emph{Selected Areas in Cryptography -- SAC 2020}, pp.~404--430.
\newblock \url{https://doi.org/10.1007/978-3-030-81652-0_16}

\bibitem{hashimoto2021}
K.~Hashimoto, S.~Katsumata, K.~Kwiatkowski, T.~Prest.
\newblock An Efficient and Generic Construction for Signal's Handshake
  (X3DH): Post-Quantum, State Leakage Secure, and Deniable.
\newblock In \emph{PKC 2021}, pp.~410--440.
\newblock \url{https://doi.org/10.1007/978-3-030-75248-4_15}

\bibitem{fips203}
National Institute of Standards and Technology.
\newblock Module-Lattice-Based Key-Encapsulation Mechanism Standard.
\newblock FIPS Publication~203, 2024.
\newblock \url{https://doi.org/10.6028/NIST.FIPS.203}

\bibitem{shor1997}
P.~W.~Shor.
\newblock Polynomial-time algorithms for prime factorization and
  discrete logarithms on a quantum computer.
\newblock \emph{SIAM J. Comput.}, 26(5):1484--1509, 1997.
\newblock \url{https://arxiv.org/abs/quant-ph/9508027}

\bibitem{cecpq2}
A.~Langley.
\newblock CECPQ2.
\newblock Google Chrome experiment, 2018.
\newblock \url{https://www.imperialviolet.org/2018/12/12/cecpq2.html}

\bibitem{stebila2020}
D.~Stebila, M.~Mosca.
\newblock Post-Quantum Key Exchange for the Internet and the Open
  Quantum Safe Project.
\newblock In \emph{Selected Areas in Cryptography -- SAC 2016}, pp.~14--37.
\newblock \url{https://doi.org/10.1007/978-3-319-69453-5_2}

\bibitem{dowling2020}
B.~Dowling, M.~Fischlin, F.~G{\"u}nther, D.~Stebila.
\newblock A Cryptographic Analysis of the TLS 1.3 Handshake Protocol.
\newblock \emph{Journal of Cryptology}, 34(4):37, 2021.
\newblock \url{https://doi.org/10.1007/s00145-021-09384-1}

\bibitem{rfc9420}
R.~Barnes, B.~Beurdouche, R.~Robert, J.~Millican, E.~Omara,
  K.~Cohn-Gordon.
\newblock The Messaging Layer Security (MLS) Protocol.
\newblock RFC~9420, IETF, 2023.
\newblock \url{https://www.rfc-editor.org/rfc/rfc9420}

\bibitem{rfc7748}
A.~Langley, M.~Hamburg, S.~Turner.
\newblock Elliptic Curves for Security.
\newblock RFC~7748, IETF, 2016.
\newblock \url{https://www.rfc-editor.org/rfc/rfc7748}

\bibitem{rfc5869}
H.~Krawczyk, P.~Eronen.
\newblock HMAC-based Extract-and-Expand Key Derivation Function (HKDF).
\newblock RFC~5869, IETF, 2010.
\newblock \url{https://www.rfc-editor.org/rfc/rfc5869}

\bibitem{krawczyk2010}
H.~Krawczyk.
\newblock Cryptographic Extraction and Key Derivation: The HKDF Scheme.
\newblock In \emph{CRYPTO 2010}, LNCS~6223, pp.~631--648.
\newblock \url{https://doi.org/10.1007/978-3-642-14623-7_34}

\bibitem{rogaway2006}
P.~Rogaway, T.~Shrimpton.
\newblock A Provable-Security Treatment of the Key-Wrap Problem.
\newblock In \emph{EUROCRYPT 2006}, LNCS~4004, pp.~373--390.
\newblock \url{https://doi.org/10.1007/11761679_23}

\bibitem{rfc8452}
S.~Gueron, A.~Langley, Y.~Lindell.
\newblock AES-GCM-SIV: Nonce Misuse-Resistant Authenticated Encryption.
\newblock RFC~8452, IETF, 2019.
\newblock \url{https://www.rfc-editor.org/rfc/rfc8452}

\bibitem{sealed-sender}
J.~Lund.
\newblock Technology preview: Sealed sender for Signal.
\newblock Signal Blog, 2018.
\newblock \url{https://signal.org/blog/sealed-sender/}

\bibitem{sphinx}
G.~Danezis, I.~Goldberg.
\newblock Sphinx: A compact and provably secure mix format.
\newblock In \emph{IEEE S\&P 2009}, pp.~269--282.
\newblock \url{https://doi.org/10.1109/SP.2009.15}

\bibitem{loopix}
A.~Piotrowska, J.~Hayes, T.~Elahi, S.~Meiser, G.~Danezis.
\newblock The Loopix Anonymity System.
\newblock In \emph{USENIX Security 2017}, pp.~1199--1216.
\newblock \url{https://www.usenix.org/conference/usenixsecurity17/technical-sessions/presentation/piotrowska}

\bibitem{acd2019}
J.~Alwen, S.~Coretti, Y.~Dodis.
\newblock The Double Ratchet: Security Notions, Proofs, and
  Modularization for the Signal Protocol.
\newblock In \emph{EUROCRYPT 2019}, LNCS~11476, pp.~129--158.
\newblock \url{https://doi.org/10.1007/978-3-030-17653-2_5}

\bibitem{bienstock2020}
A.~Bienstock, J.~Fairoze, S.~Garg, P.~Mukherjee, S.~Raghuraman.
\newblock A More Complete Analysis of the Signal Double Ratchet Algorithm.
\newblock In \emph{CRYPTO 2022}, LNCS~13507, pp.~784--813.
\newblock \url{https://doi.org/10.1007/978-3-031-15802-5_27}

\bibitem{bienstock2022}
A.~Bienstock, Y.~Dodis, P.~R{\"o}sler.
\newblock On the Price of Concurrency in Group Ratcheting Protocols.
\newblock In \emph{TCC 2020}, pp.~198--228.
\newblock \url{https://doi.org/10.1007/978-3-030-64378-2_8}

\bibitem{canetti2001}
R.~Canetti, H.~Krawczyk.
\newblock Analysis of Key-Exchange Protocols and Their Use for Building
  Secure Channels.
\newblock In \emph{EUROCRYPT 2001}, pp.~453--474.
\newblock \url{https://doi.org/10.1007/3-540-44987-6_28}

\bibitem{lamacchia2007}
B.~LaMacchia, K.~Lauter, A.~Mityagin.
\newblock Stronger Security of Authenticated Key Exchange.
\newblock In \emph{ProvSec 2007}, LNCS~4784, pp.~1--16.
\newblock \url{https://doi.org/10.1007/978-3-540-75670-5_1}

\bibitem{meier2013}
S.~Meier, B.~Schmidt, C.~Cremers, D.~Basin.
\newblock The TAMARIN Prover for the Symbolic Analysis of Security
  Protocols.
\newblock In \emph{CAV 2013}, LNCS~8044, pp.~696--701.
\newblock \url{https://doi.org/10.1007/978-3-642-39799-8_48}

\bibitem{blanchet2016}
B.~Blanchet.
\newblock Modeling and Verifying Security Protocols with the Applied
  Pi Calculus and ProVerif.
\newblock \emph{Foundations and Trends in Privacy and Security},
  1(1--2):1--135, 2016.
\newblock \url{https://bblanche.gitlabpages.inria.fr/publications/BlanchetFnTPS16.html}

\bibitem{aura-security-proof}
O.~Melnychenko.
\newblock Aura Protection Protocol: Game-based security proofs.
\newblock Technical report, 2025.
\newblock Companion technical report with full proofs of the six
  security theorems; local companion file: \path{docs/security-proof.tex}.

\end{thebibliography}

\end{document}
